\documentclass[11pt]{article}

\usepackage[a4paper,margin=2.5cm]{geometry}
\usepackage{fontspec}
\setmonofont{Consolas}
\usepackage{amsmath,amssymb,amsthm}
\usepackage{booktabs}
\usepackage{xcolor}
\usepackage{hyperref}

\newtheorem{remark}{Remark}

\newcommand{\ms}{\mathrm{ms}}
\newcommand{\any}{\mathrm{any}}
\newcommand{\free}{\mathrm{free}}
\newcommand{\ksel}{\mathrm{ksel}}
\newcommand{\ong}{\mathrm{ong}}
\newcommand{\pole}{\mathrm{pole}}
\DeclareMathOperator*{\argmax}{arg\,max}

\title{Formulations implemented in \texttt{paper\_experiment}}
\author{}
\date{}

\begin{document}
\maketitle

\noindent
\textcolor{red}{We will probably need to lower the MTBF values so that every component is
significant.}

\noindent
\textcolor{red}{Since with the current numerical values maintenance start days were
already effectively restricted to day $1$, $1+\mathrm{QS}$ or $1+2\,\mathrm{QS}$ of the
period (every repair duration is a multiple of $\mathrm{QS} = 20$), the pricing MILP of
Section~\ref{sec:pricing} (\texttt{solve\_pricing} in \texttt{fluid\_lp\_cg.jl}, and its
counterpart \texttt{solve\_pricing\_light} in \texttt{fluid\_cg\_light.jl}) now enforces
this restriction explicitly: the start-day variables $m^{\rm start}_{m,\tau}$ are only
created for $\tau$ on this grid, and the technician-capacity constraint (11) is only
imposed at the grid points instead of on every day $\tau=1,\dots,T$. Because the daily
penalty is currently uniform over the period, this is \emph{exact}, not an approximation:
any feasible schedule can be shifted onto the grid without changing its cost (only the
\emph{number} of penalised days matters, not which days they are) and without breaking
technician-capacity feasibility (shifting starts onto the grid only tightens/aligns
overlaps, it cannot create a new one). It cuts pricing-MILP solve time substantially (the
number of start-day
variables and technician-capacity constraints drops from $O(T)$ to $O(T/\mathrm{QS})$
per component). \textbf{This stops being exact if the daily penalty is made
non-uniform}: the start day would then affect cost, and this grid restriction would
have to be revisited (either dropped, or re-justified against the new penalty
structure) before trusting the resulting LP bound.}

\noindent
This note transcribes, as literally as possible, the three optimisation problems coded in
\texttt{paper\_experiment/fluid\_lp\_cg.jl} (master LP and pricing MILP of the Full Fluid LP,
solved by column generation) and \texttt{paper\_experiment/compact\_lp.jl} (Compact Fluid LP).
The formulations correspond respectively to Sections 3.1--3.2 (full fluid LP / column
generation) and Section 3.3 (compact fluid LP) of the paper, specialised to the offshore
substation instance of Section 5 and Appendix C. Index names follow the code exactly
(function and variable names are given in footnotes) so that every symbol below can be
matched to a line of Julia.

\section{Notation and data}
\label{sec:notation}

\paragraph{Poles and components.} An instance has $K$ poles $p = 1,\dots,K$
(\texttt{inst.n\_poles}). Each pole is a set of $7$ components (converter, water circuit,
fans, pumps, transformer, cooling, busbar); flattening all poles gives $M$ components
$m = 1,\dots,M$ (\texttt{fd.M}), with $\pole(m)$ the pole of $m$ and
$\mathrm{comps\_of\_pole}(p)$ the (sorted) list of components of pole $p$; we write
$m_1(p) := \mathrm{comps\_of\_pole}(p)[1]$ for the \emph{reference component} of pole $p$.

\paragraph{Component state and cost.} Component $m$ has $n_m$ degradation states
$x \in \{1,\dots,n_m\}$ (\texttt{fd.n\_states[m]}), $x=1$ being ``as new''. A \emph{mode}
$a \in \{1,2\}$ indicates whether $m$ is left alone ($a=1$) or maintained/replaced-as-new
($a=2$) during the period. The per-period cost of $(x,a)$ is
\[
  c_m(x,a) \;=\; \begin{cases}
    \delta_m(x) & a = 1 \quad (\text{expected failure penalty, \texttt{fd.δ[m][x]}}) \\
    \mathrm{cc}_m & a = 2 \quad (\text{replacement cost, \texttt{fd.cc[m]}}).
  \end{cases}
\]
Degradation follows $p^1_m(x \mid x') := P_m(x \mid x')$ (\texttt{fd.P[m][x,x']}, the
period transition matrix when $a=1$), while maintenance resets deterministically to the
as-new state: $p^2_m(x \mid x') := \mathbb 1\{x = 1\}$ (independent of $x'$).

\paragraph{$\alpha$ (\texttt{build\_fluid\_data(inst; alpha=1.0)}).} $\delta_m(x)$ above is
always $\alpha$ times the cost of leaving $m$ unmaintained this period in state $x$ --
\texttt{failure\_cost} (day-by-day expectation) or \texttt{failure\_cost\_simplified} (cost
of staying in $x$ for the whole period, Section~\ref{sec:simplified}), whichever
\texttt{inst.simplified} selects. $\alpha = 1$ (default) leaves that cost as defined by
either formula; $\alpha \ne 1$ reweights the failure penalty relative to the (always
unscaled) replacement cost $\mathrm{cc}_m$ and joint schedule cost $f(b)$, e.g. to study how
sensitive a method's policy is to that trade-off. $\alpha$ is fixed once, when
\texttt{fd} is built (\texttt{run\_experiment.jl}'s \texttt{--alpha=<value>}, default
$1.0$): every method reads it only through the already-scaled \texttt{fd.δ}, so comparing
methods at different $\alpha$ means building one \texttt{FluidData} per $\alpha$ and
passing each to the methods separately, not passing $\alpha$ to a method's own solve
function.

\paragraph{Quota.} Each pole $p$ has an annual quota $K_0(p)$ (\texttt{fd.K0[p]}) of
penalty-free maintenance days. The \emph{local state} of component $m$ at stage $t$ is the
pair $(x,k)$, $k \in \{0,\dots,K_0(\pole(m))\}$ being the quota remaining at the
\emph{start} of period $t$; the \emph{local action} is $(a,\kappa)$, $\kappa \in
\{0,\dots,k\}$ being the number of free days of the pole used \emph{during} period $t$.
Every component of a pole is required to carry the same $(k,\kappa)$ distribution
(Section 5.3 of the paper). We write $\bar K := \max_p K_0(p)$ (\texttt{maximum(fd.K0)}).
In the code, $k$ and $\kappa$ do not range over every integer in $\{0,\dots,K_0(p)\}$: they
are restricted to multiples of a fixed step $\mathrm{QS} = 20$ (\texttt{QUOTA\_STEP}), i.e.\
$k,\kappa \in \{0, \mathrm{QS}, 2\,\mathrm{QS}, \dots\} \cap \{0,\dots,K_0(p)\}$ (helper
function \texttt{qrange}). This is a coarsening of the quota state/action space, not a
change of model: it reduces the number of $(k,\kappa)$ pairs the master LP and pricing
MILP have to enumerate, at the cost of only tracking the quota balance to the nearest
$\mathrm{QS}$ days. \textbf{Every equation below that quantifies over $k \in
\{0,\dots,K_0(p)\}$ or $\kappa \in \{0,\dots,k\}$ (or the analogous ranges for
$\bar K$) should be read as implicitly restricted to this stepped grid}; we keep the
$\{0,\dots,\cdot\}$ notation throughout for readability rather than writing
\texttt{qrange}$(\cdot)$ everywhere.

\paragraph{Schedule data.} $T$ is the number of days per period (\texttt{inst.T\_period}),
$d_m$ the repair duration of $m$ (\texttt{fd.d[m]}), $\rho_m$ its technician requirement
(\texttt{fd.rho[m]}), $R$ the total available workforce (\texttt{inst.R}), and
$\mathrm{margin}$ the number of trailing days of a period during which no maintenance may
start (\texttt{inst.margin}). The last admissible start day of $m$ is
\[
  L_m \;:=\; T - d_m + 1 - \mathrm{margin} \qquad (\texttt{last\_start[m]}).
\]
If $L_m < 1$, component $m$ can never be maintained and is dropped from the schedule
variables (its $z_m$ is fixed to $0$). $\mathrm{pen}_p(\tau)$ denotes the capacity-loss
penalty coefficient of pole $p$ on day $\tau$ of the period
(\texttt{inst.poles[p].daily\_penalty[\ensuremath{\tau}]}).

\paragraph{Lead margin.} No maintenance may start during the first $\mathrm{LEAD\_MARGIN}$
days of a period either (\texttt{LEAD\_MARGIN}, \texttt{instances.jl}): it takes that many
days after a maintenance is planned before it can actually start (ordering parts, etc.).
The earliest admissible start day of every component is therefore
\[
  F \;:=\; \mathrm{LEAD\_MARGIN} + 1 \qquad (\texttt{first\_start}),
\]
the same for every component. If $L_m < F$, component $m$ can never be maintained this
period (same fallback as $L_m < 1$ above). \textbf{Every equation below that quantifies a
start day $\tau$ or a start-day variable $s_{m,\tau}$ (or $m^{\mathrm{start}}_{m,\tau}$)
over $\{1,\dots,L_m\}$ should be read as restricted to $\{F,\dots,L_m\}$}, and every sum
$\sum_{\tau=1}^{L_m}$ as $\sum_{\tau=F}^{L_m}$ -- we keep the $\{1,\dots,L_m\}$ notation
throughout for readability, exactly as for the $\mathrm{qrange}$ convention above.

\paragraph{Numerical values (\texttt{paper\_experiment/instances.jl}).} $T = 60$
(\texttt{T\_PERIOD}), $\mathrm{margin} = 0$ (\texttt{MARGIN}),
$\mathrm{LEAD\_MARGIN} = 15$ (so $F = 16$), $R = 1$
(\texttt{inst.R}), and $K_0(p) = 30$ for every pole (\texttt{FREE\_QUOTA}), with quota
step $\mathrm{QS} = 15$ (\texttt{QUOTA\_STEP}), so that $k$ and $\kappa$ only take the
values $0$, $15$ and $30$. The
technician requirement $\rho_m$ and repair duration $d_m$ are, per component (converter,
water circuit, fans, pumps, transformer, converter cooling, busbar):
\begin{center}
\begin{tabular}{lccccccc}
\toprule
component & converter & water circuit & fans & pumps & transformer & cooling & busbar \\
\midrule
$\rho_m$ & $1.0$ & $0.5$ & $0.5$ & $0.5$ & $1.0$ & $0.5$ & $0.5$ \\
$d_m$    & $30$  & $15$  & $15$  & $15$  & $15$  & $15$  & $15$  \\
\bottomrule
\end{tabular}
\end{center}
(\texttt{TECHNICIAN\_SHARE}, \texttt{REPAIR\_DURATION}). Every component but the converter
and the transformer needs half the workforce ($\rho_m=0.5$) for $15$ workdays; the converter
and the transformer need the whole workforce ($\rho_m=1$), the converter for $30$ workdays
(twice as long as every other component, including the transformer). Since two components
with $\rho_m=1$ can never overlap, and at most two with $\rho_m=0.5$ can run simultaneously,
these values make the technician-capacity constraint (11) of Section~\ref{sec:pricing} bind
tightly enough to be worth studying on its own: with $L_m$ and $F$ as above, the admissible
scheduling window has length $L_m - F + 1 = 45$ days for every component ($T - \mathrm{margin}
- \mathrm{LEAD\_MARGIN} = 60 - 0 - 15$), exactly $3/4$ of the $60$-day window of the
previous parameters ($\mathrm{margin}=0$, no lead margin), and every duration $d_m$ is also
exactly $3/4$ of its previous value ($40 \to 30$, $20 \to 15$). This is not a coincidence:
the new parameters were chosen as a uniform $3/4$ rescaling of the previous ones, so the
feasibility structure of the scheduling problem (which subsets of components fit
simultaneously in the window, for any given technician-capacity budget) is exactly
preserved -- one pole alone still admits $80$ distinct subsets of components with a
feasible schedule within one period (including the empty one), and the whole $3$-pole
region (technician capacity shared across poles) still admits $16\,508$.

\paragraph{Stages.} $t = 1,\dots,H$ index the periods of the fluid horizon
(\texttt{H}). A calendar-year reset of the quota occurs between stage $t$ and $t+1$ iff
$t$ is the sixth period of its cyclic year (function \texttt{stage\_cyc}); we simply write
``$t \!\to\! t{+}1$ resets'' for this event below.

\section{Simplified problem: no free-day quota, no within-period transition}
\label{sec:simplified}

Instances \texttt{one\_pole\_simple}, \texttt{two\_poles\_simple} and
\texttt{three\_poles\_simple} (\texttt{paper\_experiment/instances.jl},
\texttt{build\_instance(n\_poles; free\_quota = 0, name\_suffix = "\_simple")}) drop the
free-day quota entirely, while keeping exactly the same poles, components, technician
shares, repair durations and replacement costs as \texttt{one\_pole}/\texttt{two\_poles}/
\texttt{three\_poles} (Section~\ref{sec:notation}) -- so the set of realisable maintenance
\emph{schedules} $b$ (which components, on which start days, under the same
technician-capacity constraint (11)) is completely unchanged. Only the cost is simplified,
to land exactly on the additive formalism of Section 2 of the paper: a sum of per-component
local costs plus a single cost term that is a function of the joint schedule alone, with no
quota bookkeeping in between.

\paragraph{How $K_0(p) = 0$ gives this for free.} Setting the annual quota
$K_0(p) = 0$ (\texttt{free\_quota = 0}) for every pole, with every other data item
unchanged, already produces exactly this simplified problem, with \emph{no change to any
of the LP/MILP code} of Sections~\ref{sec:master}--\ref{sec:lift-project}: since
$\mathrm{qrange}(0) = 0\!:\!\mathrm{QS}\!:\!0 = \{0\}$ (Section~\ref{sec:notation}), every
pole's $(k,\kappa)$ collapses to the single value $(0,0)$ throughout, so no component can
ever have $\kappa > 0$, and \texttt{pole\_penalty}'s (\texttt{instances.jl:109--117})
\texttt{free\_day} branch is never taken. The joint cost of a schedule $b$ then reduces to
\[
  f(b) \;=\; \sum_{p=1}^{K} \sum_{\tau=1}^{T} \mathrm{pen}_p(\tau) \cdot \mathbb 1[\text{some component of pole } p \text{ is under maintenance on day } \tau],
\]
and since $\mathrm{pen}_p(\tau) \equiv 1$ (\texttt{DAILY\_PENALTY}, Section~\ref{sec:notation}),
$f(b)$ is exactly, per pole, the \emph{number of days that pole is under maintenance},
summed over poles -- the ``maintenance cost'' requested for the simplified problem. Total
period cost is then
\[
  \sum_{m=1}^{M} c_m(x_m, a_m) \;+\; f(b),
\]
with, on a simplified instance,
\[
  c_m(x,1) = \delta_m(x), \qquad c_m(x,2) = \delta_m(x) + \mathrm{cc}_m
  \qquad (\texttt{local\_cost}, \texttt{fluid\_lp\_cg.jl}).
\]
The failure cost $\delta_m(x)$ of the period spent in state $x$ is paid \emph{whatever the
mode}, also when $m$ is maintained during the period (it is not repaired instantly), and a
maintained component pays in addition its replacement cost $\mathrm{cc}_m$. (On the original
instances, $c_m(x,2) = \mathrm{cc}_m$ only, as in Section~\ref{sec:notation}.) This is
exactly the sum of per-subsystem local costs plus one joint-schedule cost term, with $b$
ranging over the same combinatorial schedules as before.

\paragraph{Running every method unchanged.} Because $K_0(p) = 0$ only trivialises the
$(k,\kappa)$ dimension already present in every formulation's variables (it does not
remove or add any variable, constraint or code path), \texttt{fluid\_cg},
\texttt{fluid\_cg\_light}, \texttt{compact}, \texttt{fluid\_cg\_order2} and
\texttt{lift\_project} all run on \texttt{*\_simple} exactly as on the original instances,
selected by instance name alone (\texttt{run\_experiment.jl <method> one\_pole\_simple
...}). Verified directly (with $\mathrm{LEAD\_MARGIN}=15$, $d_m \in \{30,15\}$,
$K_0=30$/$\mathrm{QS}=15$ as above): on \texttt{one\_pole\_simple}, all five methods
converge (\texttt{fluid\_cg}: $0.0821$, \texttt{fluid\_cg\_light}: $0.0821$,
\texttt{compact}: $0.0793$, \texttt{fluid\_cg\_order2}: $0.0821$, \texttt{lift\_project}:
$0.0896$), the same ordering (\texttt{compact} loosest, \texttt{lift\_project} tightest) as
on the quota-augmented instances.

\paragraph{No within-period transition (\texttt{instances.jl}, \texttt{Instance.simplified}).}
On top of $K_0(p) = 0$, \texttt{*\_simple} instances also set \texttt{simplified = true},
which removes the day-by-day degradation transition entirely. A component that is not
maintained this period no longer evolves day by day through \texttt{P\_daily}: it keeps
its entering state $x$ for the whole period, and only transitions between periods, through
$P^{\mathrm{period}} = (P^{\mathrm{daily}})^{T}$ (already the transition kernel $p^1_m$ used
by the fluid LP itself, Section~\ref{sec:notation}) -- exactly the assumption ``no
degradation in progress during the period'' requested for this instance. A maintained
component still resets to the as-new state the moment its repair completes, unaffected.

\emph{Failure cost.} $\delta_m(x)$ is no longer the day-by-day
expectation of \texttt{failure\_cost} (propagating $x$ through $P^{\mathrm{daily}}$ every
day, Section~\ref{sec:notation}): it is the cost of staying in state $x$, unchanged, for
the whole period,
\[
  \delta_m(x) \;=\; \alpha \cdot \Big(\sum_{\tau=1}^{T} \mathrm{pen}_p(\tau)\Big) \cdot \frac{100 - P_{\mathrm{evac}}(c,x)}{100}
  \qquad (\texttt{failure\_cost\_simplified}, \texttt{fluid\_lp\_cg.jl}),
\]
where $c$/$p$ are the component/pole of $m$ and $\alpha$ is as in Section~\ref{sec:notation}
($=1$ by default) -- $T$ repetitions of one day's capacity-loss penalty at the (frozen)
state $x$, instead of an expectation over the state's day-by-day drift.

\paragraph{The true cost is no longer computed separately.} Since the request is to treat
the simplified problem's own additive cost model as the true model (not to measure a gap
against it), \texttt{simulate\_period} (\texttt{simulation.jl}) does not run the day-by-day
\texttt{pole\_penalty} accounting at all on a \texttt{simplified} instance: the period cost
is read directly off $b = (\mathrm{start}_m)_m$ and the entering state $x$,
\[
  \Big(\sum_{m=1}^{M} \delta_m(x_m)\Big) \;+\; f(b) \quad\text{and}\quad \sum_{m\,:\,a_m = 2} \mathrm{cc}_m,
\]
split into the reported \texttt{penalty} and \texttt{replacement\_cost} columns -- together
exactly $\mathtt{lp\_approx\_cost(fd, x, col)}$ (\texttt{fluid\_lp\_cg.jl}) by construction,
not merely by empirical coincidence: on a \texttt{simplified} instance, \texttt{mean\_cost}
and \texttt{mean\_lp\_cost} of \texttt{run\_experiment.jl}'s report are therefore always
exactly equal (verified: model gap $= 0.00$ on \texttt{one\_pole\_simple}, every scenario).
This is deliberately \emph{not} extended to the original, non-simplified instances: their
report keeps tracking \texttt{mean\_cost} (true, non-additive, day-by-day simulated cost)
against \texttt{mean\_lp\_cost} (the additive proxy) separately, since the gap between the
two is exactly what those instances are for (Appendix C.4 of the paper;
Section~\ref{sec:heuristic-extraction}'s discussion of why that gap matters).

\paragraph{Data of the simplified instances.} \texttt{*\_simple} instances use the same
repair durations and replacement costs as the original instances (Section~\ref{sec:notation}),
but lower MTBFs and higher technician shares (\texttt{MTBF\_SIMPLE},
\texttt{TECHNICIAN\_SHARE\_SIMPLE}, \texttt{instances.jl}; the daily matrices are rebuilt with
the same constructors as \texttt{build\_MC\_utils.jl}, and $P^{\mathrm{period}} =
(P^{\mathrm{daily}})^{T}$):
\begin{center}
\begin{tabular}{lccccccc}
\toprule
component & converter & water circuit & fans & pumps & transformer & cooling & busbar \\
\midrule
MTBF (years), original   & $588$ & $43$ & $43$ & $43$ & $2222$ & $122$ & $4762$ \\
MTBF (years), simplified & $10$  & $4$  & $4$  & $4$  & $12$   & $6$   & $15$   \\
$\rho_m$, original       & $1.0$ & $0.5$ & $0.5$ & $0.5$ & $1.0$ & $0.5$ & $0.5$ \\
$\rho_m$, simplified     & $1.0$ & $0.3$ & $0.3$ & $0.3$ & $1.0$ & $0.3$ & $0.3$ \\
\bottomrule
\end{tabular}
\end{center}
With the paper's MTBFs, the converter, transformer and busbar essentially never fail over the
horizon, so most states are healthy and there is rarely more than one component worth
maintaining at once. Lower MTBFs put several components in degraded, maintenance-worthy
states at the same time, which is when the choice of which ones to maintain together (and
how many can fit, with $\rho_m = 0.3$ allowing up to three at once instead of two with
$\rho_m = 0.5$) matters, and the relaxations differ. They also differ from the original
instances by the absence of quota and of within-period transition, and by the cost structure
described above. The MTBF remark at the top of this note applies to the original instances
only.

\paragraph{\textcolor{red}{Parameters used for the simplified model.}}
\textcolor{red}{Every \texttt{*\_simple} instance (\texttt{one\_pole\_simple},
\texttt{two\_poles\_simple}, \texttt{three\_poles\_simple}) uses the following values.}
\begin{itemize}
  \item \textcolor{red}{Structure: $K \in \{1,2,3\}$ poles of $7$ components, $n\_\mathrm{years} = 20$
  years of $6$ periods, $T = 60$ days per period, fluid horizon $H = 6$ periods, $R = 1$
  shared by all poles.}
  \item \textcolor{red}{Scheduling: $\mathrm{margin} = 0$, $\mathrm{LEAD\_MARGIN} = 15$ (so
  $F = 16$), repair durations $d_m = 30$ (converter) and $15$ (every other component),
  $\mathrm{QS} = 15$.}
  \item \textcolor{red}{Quota: $K_0(p) = 0$ (no free day).}
  \item \textcolor{red}{Costs: $\alpha = 1$, $\mathrm{pen}_p(\tau) \equiv 1$,
  replacement costs $\mathrm{cc}_m = (5,\,0.25,\,0.125,\,0.25,\,10,\,0.5,\,0.25)$ (converter,
  water circuit, fans, pumps, transformer, cooling, busbar), $\delta_m(x) = \alpha \cdot T \cdot
  (100 - P_{\mathrm{evac}}(c,x))/100$.}
  \item \textcolor{red}{Technician shares $\rho = (1,\,0.3,\,0.3,\,0.3,\,1,\,0.3,\,0.3)$
  (the original instances have $0.5$ instead of $0.3$).}
  \item \textcolor{red}{MTBF (years) $= (10,\,4,\,4,\,4,\,12,\,6,\,15)$ (the original instances
  have $(588,\,43,\,43,\,43,\,2222,\,122,\,4762)$).}
  \item \textcolor{red}{No within-period transition (\texttt{simplified = true}), one
  transition per period through $P^{\mathrm{period}} = (P^{\mathrm{daily}})^{60}$.}
\end{itemize}
\textcolor{red}{\textbf{Caveat:} these MTBFs are not very realistic. They were lowered only
so that several components degrade noticeably over the horizon and the methods can be told
apart; they are much shorter than the paper's values (e.g.\ $10$ years instead of $588$ for
the converter, $12$ instead of $2222$ for the transformer, $15$ instead of $4762$ for the
busbar).}

\paragraph{Order-2 formulation on a simplified instance.} Since the true cost of a
simplified instance is additive, there is no joint (non-additive) capacity term for the
pairwise pseudomarginals of Section~\ref{sec:order2} to correct: $\Delta_{m_1,m_2} \equiv 0$
for every same-pole pair (\texttt{compute\_Δ\_pairs} sets it directly to the zero matrix
when \texttt{inst.simplified} holds, without calling \texttt{δ\_joint\_pair} -- that
function propagates day by day, which is only meaningful against the \emph{original}
model's own day-by-day $\delta$, not the frozen-state $\delta$ of a simplified instance).
This does \emph{not} mean \texttt{fluid\_cg\_order2} reduces to \texttt{fluid\_cg} on a
simplified instance, though: \texttt{build\_master\_order2} still builds the pairwise
pseudomarginals $p_{\mathrm{pair}}$ (their own Flow constraints and their marginal
consistency with $q$) for every same-pole pair requested -- only the objective correction
term is (trivially, since $\Delta \equiv 0$) zero. The pairwise Flow/consistency
constraints are a valid tightening of the relaxation of the marginals themselves
regardless of whether there is any objective correction to make (Proposition 9-style
argument), so the order-2 LP value is still weakly tighter than (never below) the order-1
one, verified on \texttt{one\_pole\_simple} and \texttt{two\_poles\_simple} with Gurobi's
default tolerances loosened to $10^{-9}$ (\texttt{OptimalityTol}, \texttt{FeasibilityTol}):
with the default $10^{-6}$ tolerance, the two independent column-generation runs (order-1
and order-2 explore different columns, hence different, possibly degenerate, optimal
bases) can each stop up to that relative tolerance away from the true LP optimum, which
occasionally makes the reported order-2 value come out a few $10^{-5}$ \emph{below} the
order-1 one -- a solver-tolerance artefact, not a violation of the tightening property.
\textbf{$\Delta \equiv 0$ is never used as a skip condition in \texttt{build\_master\_order2}},
for either a same-pole or a cross-pole pair: a pair with $\Delta \equiv 0$ (every same-pole
pair on a simplified instance; occasionally a same-pole pair on the original model too)
still gets its $p_{\mathrm{pair}}$ built, since the pairwise Flow/consistency constraints
tighten the relaxation of the marginals themselves regardless of whether there is any
objective correction to make.

\paragraph{Cross-pole pairs are also built, without a shared $(k,\kappa)$.} Two components
$m_1,m_2$ of \emph{different} poles do not share a single quota state to index
$p_{\mathrm{pair}}$ against the way a same-pole pair shares its one pole's $(k,\kappa)$
(Remark~\ref{rem:strengthen}) -- each pole has its own, independent quota. This does not
mean a cross-pole pair is useless, though: the pairwise Flow constraint (the product kernel
$p_{m_1}^{a_1}(x_1\mid x_1') \, p_{m_2}^{a_2}(x_2\mid x_2')$) still enforces that the JOINT
$(x_1,x_2)$ trajectory is consistent with the two components evolving independently given
their own local actions (Definition 1 of the paper) -- a genuine property of the model that
$q$'s own, per-component Flow constraints do not enforce on their own (they only pin each
component's OWN marginal, never the correlation between two components' trajectories).
\texttt{build\_master\_order2} therefore builds, for a cross-pole pair, $p_{\mathrm{pair}}$
variables $p_{x_1,x_2,a_1,a_2,t}$ \textbf{without} a $(k,\kappa)$ index, with the same
Initial condition and Flow as the same-pole case (just without any reset/quota bookkeeping,
irrelevant here), and Marginal consistency tying $p_{\mathrm{pair}}$ to $q$'s own marginal
\emph{summed over} $(k,\kappa)$ instead of a shared slice of it:
\[
  \sum_{x_2,a_2} p_{x_1,x_2,a_1,a_2,t} = \sum_{k,\kappa} q_{m_1,x_1,a_1,k,\kappa,t}
  \qquad\text{(and symmetrically for $m_2$).}
\]
$\Delta \equiv 0$ always for a cross-pole pair (the penalty decomposes additively across
poles, Section~\ref{sec:notation}), so the objective correction it contributes is always
exactly $0$ -- but, as above, this is not why it is built: the tightening comes entirely
from the Flow/consistency constraints. Verified (same $10^{-9}$-tolerance check as above)
on \texttt{two\_poles\_simple}: a cross-pole pair's order-2 LP value is weakly tighter than
(never below) the order-1 one, same as a same-pole pair.

The \texttt{max\_pairs} ranking of \texttt{compute\_Δ\_pairs} ranks every same-pole pair
before every cross-pole pair, by decreasing impact $\max|\Delta|$ only as a tiebreak within
each group: on the original model that tiebreak alone already puts every same-pole pair
before every cross-pole one (cross-pole impact is always $0$, same-pole impact generally is
not), but on a simplified instance every pair (same-pole or not) has impact $0$, so the
explicit same-pole-first key is what keeps same-pole pairs prioritised there too -- cross-pole
pairs are only ever added once \texttt{max\_pairs} exceeds the number of same-pole pairs
available, filling the remaining budget rather than being skipped or wasted.
On the original instances $\Delta_{m_1,m_2}$ for a same-pole pair is unchanged (day-by-day
joint cost of \texttt{$\delta$\_joint\_pair} minus the two individual $\delta$'s).

\section{Full Fluid LP of Section 3.2 of the paper (\texttt{fluid\_cg})}
\label{sec:paper-fluid}

\textbf{The method \texttt{fluid\_cg} of \texttt{run\_experiment.jl} is now exactly the Full
fluid LP of Section 3.2 of \texttt{paper.tex}} (\texttt{fluid\_cg\_paper.jl}), with its fluid
policy of Section 3.4. The only departure from the paper is the discretisation of the quota
in steps of $\mathrm{QS}$ days (Section~\ref{sec:notation}), which is kept. The former
implementation, described in Sections~\ref{sec:master} and~\ref{sec:pricing}, is still
available as \texttt{fluid\_cg\_legacy}. The methods \texttt{compact}, \texttt{fluid\_cg\_order2}
and \texttt{lift\_project} now also follow Sections 3.3, 4.2 and 4.3 of the paper (see below),
their former implementations being available as \texttt{compact\_legacy},
\texttt{fluid\_cg\_order2\_legacy} and \texttt{lift\_project\_legacy}.

\paragraph{True cost and simulator.} On the original instances (with quota), the simulator
(\texttt{simulate\_period}, \texttt{simulation.jl}) now evaluates the true cost $\mathrm{Tc}$
of Section 2.2 of the paper (\texttt{true\_period\_cost}, \texttt{fluid\_lp\_cg.jl}). There is
no degradation during the period. The state of component $m$ on day $\tau$ is $1$ if its repair
is complete on day $\tau$ and its entering state otherwise. The capacity loss of pole $p$ on day
$\tau$ is $1$ if some component of $p$ is under repair, and the maximum of the capacity losses
$\ell_m = (100 - P_{\mathrm{evac}})/100$ of its components otherwise. The penalty is $\pi$ times
this loss on every day that is not declared penalty-free, the declared days exempting every
loss, due to maintenance or to degradation. A maintained component returns to the new state,
and every other component transitions once, through $P^{\mathrm{period}}$, at the end of the
period. The cost recorded for a trajectory (\texttt{total\_cost} in the result files) is this
true cost, not its separable approximation (reported separately as \texttt{lp\_cost}). The
result files now carry the suffix \texttt{\_paper} (\texttt{MODEL\_VERSION}), so that the
results and caches computed with the former cost model are never resumed. The simplified
instances keep their own simulator (Section~\ref{sec:simplified}).

\paragraph{Master LP.} For the observed state $(x_0, k_0)$ and a fluid horizon $H$, the master
LP (\texttt{build\_master\_paper}) has the variables $\mu^1_b$, $b \in \mathcal B'_1 \subseteq
\mathcal B(x_0)$, with the true cost $c_\otimes(x_0, b) = \mathrm{Tc}$ above as objective
coefficient, the variables $q^m_{(n,k),(\mathrm{mode},\kappa),t}$ for $t \in [2:H]$ and
$\kappa \le k$ (local state-action constraint), with the separable cost of the paper
\[
  \tilde c_m\big((n,k),(\mathrm{mode},\kappa)\big) = \mathbb 1_{\{\mathrm{mode} = 2\}}\,\mathrm{cc}_m
  + (T - \kappa)\,\pi\,\ell_m(n)
  \qquad (\texttt{paper\_local\_cost}),
\]
and the variables $\mu_{b,t}$, $b \in \mathcal B'_t \subseteq \mathcal B$, $t \in [2:H]$, with the
joint action cost $f(b)$ (\texttt{column\_penalty}). The constraints are $\sum_b \mu^1_b = 1$,
the initial condition
\[
  \sum_{\mathrm{mode},\kappa} q^m_{(n,k),(\mathrm{mode},\kappa),2}
  = \sum_{b \in \mathcal B'_1} p_m(n \mid n_{0,m}, \delta_m(b))\,
    \mathbb 1_{\{k = \phi_1(k_{0,p(m)}, \kappa_{p(m)}(b))\}}\, \mu^1_b,
\]
the local flow constraints for $t \in [3:H]$ (with the yearly reset of the quota), the
consistency constraints $\sum_{n,k} q^m_{(n,k),(\mathrm{mode},\kappa),t} = \sum_{b \mid A_m(b) =
(\mathrm{mode},\kappa)} \mu_{b,t}$ for $t \in [2:H]$, and the quota consistency constraints
between the components of a pole for $t \in [2:H]$.

\paragraph{Pricing problems.} At $t \ge 2$, the pricing problem is \texttt{solve\_pricing}
(Section~\ref{sec:pricing}), whose objective $f(b) + \sum_m \lambda_t^{m,A_m(b)}$ is the reduced
cost of the paper. At $t = 1$, \texttt{solve\_pricing\_first} minimises $c_\otimes(x_0, b) +
\sum_m \sum_x \nu^m_x\, p_m(x \mid x_{0,m}, A_m(b)) - \eta$ over $\mathcal B(x_0)$, i.e.\ with
the quota constraint $\kappa_p \le k_{0,p}$. The true cost is linearised with continuous
auxiliary variables, $L_{p,\tau} \ge \ell_m(n_{0,m}) - (\ell_m(n_{0,m}) - \ell_m(1))\,
\mathrm{done}_{m,\tau}$ for every $m$ of $p$ (maximum of the losses, $\mathrm{done}_{m,\tau}$
being $1$ iff the repair of $m$ is complete on day $\tau$), $u_{p,\tau} \ge \mathrm{any}_{p,\tau}
- \mathrm{free}_{p,\tau}$ (maintenance day not declared penalty-free) and $v_{p,\tau} \ge
L_{p,\tau} - \mathrm{any}_{p,\tau} - \mathrm{free}_{p,\tau}$ (day without maintenance not declared
penalty-free), with objective $\sum_m \mathrm{cc}_m z_m + \sum_{p,\tau} \pi_\tau (u_{p,\tau} +
v_{p,\tau})$. These variables only appear in the objective, with nonnegative coefficients, so
they take the values they represent at the optimum.

\paragraph{Policy.} The applied joint action is the column of largest weight $\mu^{1\star}_b$
(Section 3.4 of the paper). The shared heuristic of Section~\ref{sec:heuristic-extraction}
(\texttt{--no-heuristic}) does not apply to \texttt{fluid\_cg}. The columns found for a state
are kept, per $(t_0, H)$, to initialise the column generation of the next states.

\paragraph{Order 2 (\texttt{fluid\_cg\_order2}, Section 4.2 of the paper).} With
\texttt{pairs = same\_pole\_pairs(fd)}, the master LP also has, for every pair $G = (m_1, m_2)$ of
components of a same pole and every $t \in [2:H]$, the variables
$q^G_{(n_1,n_2,k),(d_1,d_2,\kappa),t}$, indexed by the single quota $k$ and number of declared
days $\kappa \le k$ of the pole (which enforces the constraints (eq:order-k-admissible) and
(eq:order-k-copies)), with the cost of order 2 $-(T - \kappa)\,\pi\,\min(\ell_{m_1}(n_1),
\ell_{m_2}(n_2))$ (\texttt{paper\_pair\_cost}). Its constraints are the initial condition of
order 2 at $t = 2$ (in the $\mu^1_b$), the flow of order 2 for $t \in [3:H]$ with the product
kernel $p_{m_1} p_{m_2}$, the consistency of order 2 $\sum_{x_G} q^G_{x_G,a_G,t} = \sum_{b \mid
A_G(b) = a_G} \mu_{b,t}$, and the marginalisation onto each component of the pair. Only the
pairs of a same pole are used, not every group of $\mathcal G_2$: for several poles, the pairs
of components of different poles would carry the quotas of two poles and multiply the size of
the LP, while the cost of order 2 only concerns the pairs of a same pole. This family of groups
is closed under taking subsets, so that the LP is still a relaxation (Proposition 9 applies
unchanged). The pricing problems include the duals of the constraints of order 2, through the
indicators $\mathbb 1\{A_G(b) = a_G\}$, linearised by $w = z_{m_1} \wedge z_{m_2} \wedge
\ksel_{p,\kappa}$. The transition probabilities of the pairs range down to $10^{-16}$, and the
master LP is solved without presolve (\texttt{Presolve = 0}), which avoids the numerical trouble
of Gurobi (quad precision, Markowitz tolerance) observed with presolve.

\paragraph{Lift-and-project (\texttt{lift\_project}, Section 4.3 of the paper).} With
\texttt{cuts = true}, the cuts (C) of Proposition 11 are generalised to the local actions, for the
pairs $m_1 \ne m_2$ of components of a same pole, the local states $x_i = (n_i, k)$ and the local
actions $a_i = (d_i, \kappa)$ with the same quota $k$ and number of declared days $\kappa$ (the
joint probability of different quotas of a same pole being $0$):
\[
  q^{m_1}_{x_1,a_1,t} + q^{m_2}_{x_2,a_2,t} \;\le\;
  \sum_{b \mid A_{m_1}(b) = a_1 \text{ or } A_{m_2}(b) = a_2} \mu_{b,t} \;+\; \tau^+,
\]
with $\tau^+ = \sum p_{m_1}(n_1 \mid n_1', d_1')\, p_{m_2}(n_2 \mid n_2', d_2')\,
\min\big(q^{m_1}_{(n_1',k'),(d_1',\kappa'),t-1}, q^{m_2}_{(n_2',k'),(d_2',\kappa'),t-1}\big)$, the sum
being over $(d_1', d_2')$, $(n_1', n_2')$ and the predecessor quotas $(k', \kappa')$ of $k$, for
$t \ge 3$, and with the exact joint distribution $\sum_b \mu^1_b\, p_{m_1}(n_1 \mid n_{0,m_1},
\delta_{m_1}(b))\, p_{m_2}(n_2 \mid n_{0,m_2}, \delta_{m_2}(b))\, \mathbb 1\{k = \phi_1(k_0,
\kappa(b))\}$ of the two local states for $t = 2$, since the state at $t = 1$ is fixed. Each
$\min$ is represented by a variable $z \le q_1$, $z \le q_2$. The cuts are generated by a
cutting-plane loop (at most \texttt{max\_cuts\_per\_round} most violated cuts per round, then
column generation again). The columns added after a cut enter its right-hand side
(\texttt{add\_column\_paper!}, \texttt{add\_first\_column!}), and the duals of the cuts enter
the pricing problems, so that the column generation stays exact. This corrects the former
implementation (\texttt{lift\_project\_legacy}), whose cuts ignored the columns added after them
and whose pricing problems ignored the duals of the cuts.

\paragraph{Compact fluid LP (\texttt{compact}, Section 3.3 of the paper).}
\texttt{compact\_paper.jl} builds the LP (eq:compact-fluid) for the observed state. The variables
$q^m_{x,a,t}$ are defined for every $t \in [1:H]$, with the initial condition at $t = 1$, the
local flow, the quota consistency, and the cost $\tilde c_m$ at $t \ge 2$ only. For every $t$,
the variables of the MILP formulation of $\mathcal B$ (start days on the grid, penalty-free days,
indicators of the number of declared days) are relaxed into $[0,1]$. At $t = 1$, they also
satisfy the quota constraint $\kappa \le k_0$, and carry the auxiliary variables $L$, $u$, $v$ of
the true cost, the cost of the first time step being $c_\otimes(x_0, \bar b_1) = \sum_m
\mathrm{cc}_m z_{m,1} + \sum \pi (u + v)$. At $t \ge 2$, the cost is $f(\bar b_t) = \sum \pi
(\mathrm{any} - w)$. The consistency constraints are $\sum q^m_{x,(2,\kappa),t} = z_{m,t}$ and
$\sum \kappa\, q^m_{x,(d,\kappa),t} = \sum_\tau \free_{p(m),\tau,t}$, for every $t$. The policy
maintains the components with $z_{m,1}^\star > 1/2$ (\texttt{schedule\_paper}): each one starts
on the earliest day of the grid compatible with the technician capacity, the most degraded first,
and the declared penalty-free days of each pole are its days of largest loss, maintenance days
first, their number being the largest multiple of $\mathrm{QS}$ at most the remaining quota and at
most the number of days with a positive loss.

\section{Pricing problems of \texttt{fluid\_cg} by blocks and by dynamic programming}
\label{sec:pricing-dp}

On \texttt{three\_poles} with $\mathrm{QS} = 15$, the pricing problems of
Section~\ref{sec:paper-fluid} took about $90\%$ of the computing time of the column generation.
Each MILP took about $1.4$ s and explored about $665$ nodes on average, because the declared days
were chosen day by day, although the days of a block of $15$ days are interchangeable. This
section describes two faster formulations of the same pricing problems. Both are exact. The first
one (blocks) is the default. The second one (dynamic programming) is used with
\texttt{solve\_fluid\_cg\_paper(...; pricing = :dp)}.

\subsection*{Notation}

We write $\mathcal M^p$ for the components of pole $p$ (\texttt{comps\_of\_pole}),
$\mathcal G_m = \{F, F + \mathrm{QS}, \ldots\} \cap [1 : T - d_m + 1]$ for the grid of start days
of component $m$, with $F = \texttt{FIRST\_START\_DAY}$ (\texttt{grid\_starts}), and $\mathcal C =
\{F, F + \mathrm{QS}, \ldots\} \cap [1:T]$ for the \emph{check days}, on which the technician
capacity is imposed, as in the MILP (\texttt{check\_days}). With $T = 60$, $F = 16$ and
$\mathrm{QS} = 15$, $\mathcal G_m = \{16, 31, 46\}$ for the repairs of $15$ days, $\mathcal G_m =
\{16, 31\}$ for the converter ($30$ days), and $\mathcal C = \{16, 31, 46\}$. The pricing problem
at $t \ge 2$ minimises $f(b) + \sum_m \lambda_t^{m, A_m(b)}$ over $\mathcal B$, and the first-step
pricing problem minimises $c_\otimes(x_0, b) + \sum_m \sum_n \nu^m_{(n, \phi_1(k_{0,p(m)},
\kappa_{p(m)}(b)))}\, p_m(n \mid n_{0,m}, \delta_m(b)) - \eta$ over $\mathcal B(x_0)$.

\subsection*{MILP by blocks (default)}

\texttt{build\_pricing\_core} splits the period into blocks of $\mathrm{blk}$ days, and indexes
the declared days, the maintenance indicators and the losses by block instead of by day. The
variable $\free_{p,j} \in \{0,1\}$ is $1$ iff the $\mathrm{blk}$ days of block $j$ are declared
penalty-free on pole $p$, the quota constraint becomes $\mathrm{blk} \sum_j \free_{p,j} = \sum_\kappa
\kappa\, \ksel_{p,\kappa}$, and the penalty of a block is the sum of the daily penalties of its
days. \texttt{pricing\_block\_len} sets $\mathrm{blk} = \mathrm{QS}$ when this is exact, and
$\mathrm{blk} = 1$ (the formulation by day) otherwise. The conditions are the following.
\begin{itemize}
  \item $T$, $F - 1$, every repair duration $d_m$ and every annual quota $K_0(p)$ are multiples of
        $\mathrm{QS}$. The repairs then cover whole blocks, so that the maintenance indicator and the
        capacity loss of a pole are constant on every block, and every admissible number of
        declared days $\kappa$ is a multiple of $\mathrm{QS}$.
  \item The daily penalty of every pole is constant on every block.
\end{itemize}
The days of a block then have the same cost, and declaring $\kappa$ days among them in any way
costs the same as declaring $\kappa / \mathrm{QS}$ whole blocks, chosen among the most expensive
ones. The optimal value is therefore unchanged, and the many equivalent solutions of the
formulation by day disappear. The column is rebuilt with every day of the declared blocks.

\subsection*{Dynamic programming over the poles (\texttt{pricing\_dp.jl})}

\paragraph{Decomposition.} The only constraint that couples the poles is the technician capacity.
Every other term of the reduced cost only depends on the repairs and on the declared days of a
single pole. Indeed $f(b)$ and $c_\otimes(x_0, b)$ are sums over the poles, the capacity loss of a
pole being a maximum over its own components, and the dual terms only depend on the local actions
$A_m(b) = (\delta_m(b), \kappa_{p(m)}(b))$, that is, on the repairs and on the number of declared
days of the pole of $m$.

A \emph{local schedule} of pole $p$ is a vector $\sigma = (s_m)_{m \in \mathcal M^p}$ of start days,
$s_m \in \{0\} \cup \mathcal G_m$, $s_m = 0$ meaning that $m$ is not maintained. Its \emph{load} on
the check day $c \in \mathcal C$ is
\[
  u_\sigma(c) = \sum_{m \in \mathcal M^p \mid s_m > 0,\ s_m \le c \le s_m + d_m - 1} \rho_m .
\]
We write $\mathcal S_p$ for the local schedules of $p$ such that $u_\sigma(c) \le R$ for every $c
\in \mathcal C$ (\texttt{pole\_schedules}). On \texttt{three\_poles}, $|\mathcal S_p| = 1163$.
A joint action is a local schedule $\sigma_p \in \mathcal S_p$ and a set of declared days
$\mathrm{Free}^p$ for every pole $p$, such that $\sum_p u_{\sigma_p}(c) \le R$ for every $c \in
\mathcal C$, and $|\mathrm{Free}^p| \in \mathcal K_p$, with $\mathcal K_p = \{0, \mathrm{QS},
\ldots, K_0(p)\}$ at $t \ge 2$ and $\mathcal K_p = \{0, \mathrm{QS}, \ldots, k_{0,p}\}$ at $t =
1$ (\texttt{qrange}).

\paragraph{Cost of a pole for a given number of declared days.} For a local schedule $\sigma$,
let $g_\sigma(\tau)$ be the penalty of the pole on day $\tau$ if this day is not declared
(\texttt{pole\_day\_costs}). At $t \ge 2$, $g_\sigma(\tau) = \pi_\tau$ if some component of $p$ is
under repair on day $\tau$, and $0$ otherwise. At $t = 1$, $g_\sigma(\tau) = \pi_\tau$ if some
component is under repair, and $\pi_\tau \max_{m \in \mathcal M^p} \ell_m(\tilde n_{m,\tau})$
otherwise, with $\tilde n_{m,\tau} = 1$ if the repair of $m$ is complete on day $\tau$ and
$n_{0,m}$ otherwise. For a given number $\kappa$ of declared days, the best declared days are the
$\kappa$ days of largest $g_\sigma(\tau)$, so that the minimum penalty is
\[
  \sum_{\tau = 1}^{T} g_\sigma(\tau) - G_\sigma(\kappa),
  \qquad G_\sigma(\kappa) = \text{sum of the } \kappa \text{ largest values of } g_\sigma .
\]
The value of pole $p$ with the local schedule $\sigma$ and $\kappa$ declared days is then
\[
  V_p(\sigma, \kappa) = r_\sigma + \sum_{\tau} g_\sigma(\tau) - G_\sigma(\kappa)
  + \sum_{m \in \mathcal M^p} D_m\big(\mathbb 1_{\{s_m > 0\}}, \kappa\big),
\]
with $r_\sigma = \sum_{m \mid s_m > 0} \mathrm{cc}_m$ at $t = 1$ and $r_\sigma = 0$ at $t \ge 2$,
and the dual terms $D_m(\delta, \kappa) = \lambda_t^{m, (\delta, \kappa)}$ at $t \ge 2$ and
$D_m(\delta, \kappa) = \sum_n \nu^m_{(n, \phi_1(k_{0,p}, \kappa))}\, p_m(n \mid n_{0,m}, \delta)$ at
$t = 1$. The best value of the schedule is $W_p(\sigma) = \min_{\kappa \in \mathcal K_p} V_p(\sigma,
\kappa)$ (\texttt{best\_kappa}). The quantities $r_\sigma + \sum_\tau g_\sigma(\tau)$ and
$G_\sigma(\cdot)$ do not depend on the duals. They are computed once per pole and per degradation
states of its components (\texttt{schedule\_costs}, cached), so that only the dual terms are
recomputed at every call.

\paragraph{Dynamic programming.} The pricing problem is
\[
  \min \Big\{ \sum_{p = 1}^{K} W_p(\sigma_p) \ \Big|\ \sigma_p \in \mathcal S_p,\ \sum_{p} u_{\sigma_p}(c)
  \le R \ \ \forall c \in \mathcal C \Big\}
\]
(minus $\eta$ at $t = 1$). It is solved by dynamic programming over the poles
(\texttt{solve\_pricing\_dp}), the state being the cumulated load $u \in \mathbb R^{\mathcal C}$ of
the poles already processed,
\[
  F_0(0) = 0, \qquad
  F_p(u) = \min_{\sigma \in \mathcal S_p,\ u' + u_\sigma = u} \big( F_{p-1}(u') + W_p(\sigma) \big),
  \qquad u \le R,
\]
and the optimal value is $\min_u F_K(u)$. The loads $\rho_m \in \{0.5, 1\}$ and the capacity $R = 1$
give $u(c) \in \{0, 0.5, 1\}$, hence at most $3^{|\mathcal C|} = 27$ states. Two reductions keep the
recursion small. First, among the schedules of $\mathcal S_p$ with the same load vector, only the
one of smallest $W_p$ can be part of an optimal solution, so that only this one is kept
(dominance). Second, the mobilization cost, which is paid once as soon as some component is
maintained, is handled by adding to the state a boolean ``some component is maintained''. The
optimal column is rebuilt from the choices stored along the recursion. The best columns of the
other final states with a negative reduced cost are returned as well (at most $\texttt{n\_pool}
- 1$ of them), like the solution pool of the MILP.

\paragraph{Exactness.} The dynamic programming solves the same problem as the MILP. The technician
capacity is imposed on the same check days, the start days are on the same grid, the declared days
are chosen optimally by the greedy rule for every $\kappa$, every $\kappa$ of $\mathcal K_p$ is
enumerated, and the reduced cost is a sum over the poles of terms that only depend on their local
schedule and on their number of declared days. The declared days are not restricted to blocks, so
that the dynamic programming also applies when $\kappa$ is not a multiple of $\mathrm{QS}$.

\paragraph{Why the start days must be discretised, and not the quota.} The size of the state space
is exponential in the number $|\mathcal C|$ of check days. With the grid of $15$ days,
$|\mathcal C| = 3$ and there are $27$ states. With start days on every day, the capacity would have
to be checked on $45$ days and the dynamic programming would be intractable. The quota, on the
contrary, only enters through $\mathcal K_p$, which is enumerated pole by pole. Counting the quota
by steps of one day multiplies the cost of \texttt{best\_kappa} by about $10$, but leaves the
recursion unchanged.

\paragraph{Limits.} The dynamic programming handles the duals of the Joint fluid LP ($\eta$, $\nu$,
$\lambda$), but not those of the order-2 constraints and of the lift-and-project cuts
(\texttt{solve\_pricing\_first\_dp} and \texttt{solve\_pricing\_t\_dp} raise an error if
\texttt{pairs} or \texttt{cuts} are present). Their dual terms only involve pairs of components of
a same pole, so that they could be added to $V_p$ without changing the recursion. Like the MILP,
it allows declared days among the first $L$ days of the period.

\subsection*{Numerical comparison}

On \texttt{three\_poles} ($M = 21$, $\mathrm{QS} = 15$, \texttt{MTBF\_SCALE} $= 0.5$, $H = 6$, $2$
threads), along the column generation of the MILP by blocks, the three pricing methods were solved
on the same duals ($72$ pricing problems per initial state). The reduced cost of every column
returned was recomputed independently, and its feasibility checked (technician capacity on every
day, grid of start days, quota). No error was found. The largest difference between the optimal
values is $5 \cdot 10^{-5}$ between the two MILPs (the default \texttt{MIPGap} of Gurobi), and
$5 \cdot 10^{-7}$ between the MILP by day and the dynamic programming.
\begin{center}
\begin{tabular}{lccc}
  \toprule
  Time per pricing problem & MILP by day & MILP by blocks & Dynamic programming \\
  \midrule
  new components & $0.66$ s & $0.11$ s & $0.022$ s \\
  degraded components & $0.45$ s & $0.09$ s & $0.023$ s \\
  \bottomrule
\end{tabular}
\end{center}
For a whole column generation from the state with new components, without pool of columns, the
total time is $206$ s with the MILP by day, $42$ s with the MILP by blocks and $15$ s with the
dynamic programming, of which $4$ s of pricing and $11$ s of master LP. The master LP is then the
bottleneck. From the degraded state, the column generation reaches \texttt{max\_iter} $= 50$
without converging, whatever the pricing method.

\begin{remark}[Accuracy of the master LP]
Two converged column generations, one with each pricing method, gave $0.271726$ and $0.271909$.
Neither pricing method finds a column of negative reduced cost on the final duals of the other run,
and with the tolerances \texttt{FeasibilityTol} $=$ \texttt{OptimalityTol} $= 10^{-9}$ on the master
LP, both runs give $0.271953$. With the default tolerances of Gurobi ($10^{-6}$ on several thousand
variables), the values of the fluid LP are thus only accurate to about $0.1\%$. This does not affect
the policy, which only uses the column of largest weight $\mu^{1\star}_b$.
\end{remark}

\section{Master LP of the Full Fluid LP (column generation)}
\label{sec:master}

Coded in \texttt{fluid\_lp\_cg.jl}, function \texttt{build\_master} (variables
\texttt{q}) together with \texttt{add\_column!} (variables $\mu$, added incrementally by
column generation). This is Eq.~(4) of the paper (Section 3.1), together with its
consistency-to-columns constraint of Section 3.2 (Eq.~(4c)/(5)).

\begin{remark}[Strengthening the intra-pole link]
\label{rem:strengthen}
A first version of this master LP tied the components of a pole together only through an
\emph{intra-pole quota-copy} constraint, requiring every component of a pole to have the
same marginal \emph{over $(k,\kappa)$ only} as a reference component (i.e.\ the same
probability of being at quota $k$ and using $\kappa$ free days, regardless of which joint
action $b$ produced that $\kappa$). This is strictly weaker than requiring the components
to agree on the full joint distribution of $(k,b)$. We now enforce the stronger property
directly: every component of a pole is tied to the very \emph{same} variable
$\mu^{(p)}_{k,b,t}$ below, so that ``the pole's quota is $k$ and the chosen action is $b$''
is a single shared random event for all of its components, not just a matching marginal.
The former quota-copy constraint is then implied and dropped from the master (it is kept,
unchanged, in the Compact Fluid LP of Section~\ref{sec:compact}, which has no notion of a
discrete action $b$ to refine against).
\end{remark}

\subsection*{Variables}
\begin{align*}
  q_{m,x,a,k,\kappa,t} &\ge 0
  &&\forall m \in [M],\ x \in [n_m],\ a \in \{1,2\},\
     0 \le \kappa \le k \le K_0(\pole(m)),\ t \in [H], \\
  \mu_{b,t} &\ge 0
  &&\forall b \in \mathcal B' \subseteq \mathcal B,\ t \in [H], \\
  \mu^{(p)}_{k,b,t} &\ge 0
  &&\forall p \in [K],\ b \in \mathcal B',\ k \in \{\kappa_p(b),\dots,K_0(p)\},\ t \in [H],
\end{align*}
where $\mathcal B'$ is the restricted set of columns generated so far (each column
$b = (\mathrm{start}, \mathrm{free})$ fixes a start day $\mathrm{start}_m \in \{0\} \cup
\{1,\dots,L_m\}$ for every component, $0$ meaning ``not maintained'', and a set of free
days $\mathrm{free}_p \subseteq \{1,\dots,T\}$ for every pole; \texttt{Column} struct).
Write $A_m(b) := 2$ if $\mathrm{start}_m > 0$, else $1$ (\texttt{column\_mode}), and
$\kappa_p(b) := |\mathrm{free}_p|$ (\texttt{column\_kappa}). The new variable
$\mu^{(p)}_{k,b,t}$ (only created for $k \ge \kappa_p(b)$, i.e.\ for the quota states at
which using $b$ is feasible on pole $p$) is the joint probability that pole $p$'s quota is
exactly $k$ \emph{and} the chosen action is $b$, at stage $t$; it does not depend on which
component of the pole we consider (Remark~\ref{rem:strengthen}).

\subsection*{Objective}
\[
  \min \; \sum_{m,x,a,k,\kappa,t} c_m(x,a)\, q_{m,x,a,k,\kappa,t}
  \;+\; \sum_{b \in \mathcal B',\,t} f(b)\, \mu_{b,t},
\]
where $f(b)$ is the joint (schedule) cost of column $b$ (\texttt{column\_penalty}):
\[
  f(b) \;=\; \sum_{p=1}^K \; \sum_{\tau=1}^T \mathrm{pen}_p(\tau)\cdot
    \mathbb 1\bigl\{\, \tau \notin \mathrm{free}_p \ \text{and}\ \exists\, m \in
    \mathrm{comps\_of\_pole}(p):\ \mathrm{start}_m \le \tau \le \mathrm{start}_m+d_m-1 \,\bigr\}.
\]

\subsection*{Constraints}

\textbf{Initial condition} ($x_0, k_0$ the current observed state):
\[
  \sum_{a=1}^{2}\ \sum_{\kappa=0}^{k_0(\pole(m))} q_{m,x,a,k_0(\pole(m)),\kappa,1}
  \;=\; \mathbb 1\{x = x_{0,m}\}
  \qquad \forall m \in [M],\ x \in [n_m]. \tag{Init}
\]

\textbf{Flow, no reset between $t{-}1$ and $t$} ($t \ge 2$):
\[
  \sum_{a=1}^{2}\ \sum_{\kappa=0}^{k_n} q_{m,x,a,k_n,\kappa,t}
  \;=\;
  \sum_{x' \in [n_m]}\ \sum_{a=1}^{2}\ \sum_{\kappa'=0}^{K_0(\pole(m))-k_n}
    p^a_m(x \mid x')\; q_{m,x',a,\,k_n+\kappa',\,\kappa',\,t-1}
  \tag{Flow}
\]
for all $m, x \in [n_m], k_n \in \{0,\dots,K_0(\pole(m))\}$.

\textbf{Flow, reset between $t{-}1$ and $t$} ($t \ge 2$):
\[
  \sum_{a=1}^{2}\ \sum_{\kappa=0}^{k_n} q_{m,x,a,k_n,\kappa,t} \;=\; 0
  \qquad \text{if } k_n \ne K_0(\pole(m)),
\]
\[
  \sum_{a=1}^{2}\ \sum_{\kappa=0}^{K_0(\pole(m))} q_{m,x,a,K_0(\pole(m)),\kappa,t}
  \;=\;
  \sum_{x'}\ \sum_{a=1}^{2}\ \sum_{k=0}^{K_0(\pole(m))}\ \sum_{\kappa=0}^{k}
    p^a_m(x \mid x')\; q_{m,x',a,k,\kappa,t-1}
  \qquad \text{if } k_n = K_0(\pole(m)). \tag{Flow-reset}
\]

\textbf{Consistency with the joint action, at a specific quota $k$} (this is where the
columns enter the master; $\mathrm{cons}[m,a,k,\kappa,t]$ in the code):
\[
  \sum_{x \in [n_m]}\ q_{m,x,a,k,\kappa,t}
  \;-\;
  \sum_{\substack{b \in \mathcal B':\\ A_m(b)=a,\ \kappa_{\pole(m)}(b)=\kappa}} \mu^{(\pole(m))}_{k,b,t}
  \;=\; 0
  \tag{Cons}
\]
for all $m \in [M]$, $a \in \{1,2\}$, $0 \le \kappa \le k \le K_0(\pole(m))$, $t \in [H]$.
Unlike a first version, the right-hand side is \emph{not} aggregated over $k \ge \kappa$
any more: it isolates the exact quota level $k$, and --- crucially --- uses the same
variable $\mu^{(\pole(m))}_{k,b,t}$ for \emph{every} component $m$ of the pole
(Remark~\ref{rem:strengthen}), which is what makes the $(k,b)$ joint marginal identical
across components rather than merely their $(k,\kappa)$ marginals.

\textbf{Disaggregation} (ties the new variables back to $\mu_{b,t}$; created together with
$b$ itself in \texttt{add\_column!}, so it never needs a $\mathcal B'$-independent
statement):
\[
  \sum_{k=\kappa_p(b)}^{K_0(p)} \mu^{(p)}_{k,b,t} \;=\; \mu_{b,t}
  \qquad \forall p \in [K],\ b \in \mathcal B',\ t \in [H]. \tag{Disagg}
\]

Every new column $b$ is added (\texttt{add\_column!}) by, for every $t$ and every pole
$p$: creating $\mu_{b,t}$ (objective coefficient $f(b)$) and, for every
$k \in \{\kappa_p(b),\dots,K_0(p)\}$, a variable $\mu^{(p)}_{k,b,t}$ entering
$(\mathrm{Disagg})_{p,b,t}$ with coefficient $+1$ and $(\mathrm{Cons})$ at
$\bigl(m, A_m(b), k, \kappa_p(b)\bigr)$ with coefficient $-1$, for \emph{every}
$m \in \mathrm{comps\_of\_pole}(p)$ (including the components $b$ does not maintain,
which then get a $-1$ in their own $a=1$ row).

\textbf{Non-negativity:} $q,\mu,\mu^{(p)} \ge 0$.

\begin{remark}[Why $(\mathrm{Cons})$ is needed at $a=1$ too]
\label{rem:sum-to-one}
$(\mathrm{Cons})$ is imposed for \emph{both} $a=1$ (not maintained) and $a=2$ (maintained)
--- not only for $a=2$ --- and this is what guarantees $\sum_{b \in \mathcal B'} \mu_{b,t}=1$
for every $t$, as it must be since the $\mu_{b,t}$ are meant to be a probability
distribution over $\mathcal B'$. Fix $m \in [M]$ and sum $(\mathrm{Cons})$ over
$a \in \{1,2\}$ and every valid $(k,\kappa)$:
\[
  \underbrace{\sum_{a=1}^2\ \sum_{0 \le \kappa \le k \le K_0(\pole(m))}\ \sum_{x} q_{m,x,a,k,\kappa,t}}_{=\ 1
  \text{ (total probability mass of component } m \text{, by (Init)/(Flow)/(Flow-reset))}}
  \;=\;
  \sum_{a,k,\kappa}\ \sum_{\substack{b \in \mathcal B':\\ A_m(b)=a,\ \kappa_{\pole(m)}(b)=\kappa}} \mu^{(\pole(m))}_{k,b,t}.
\]
Every $b \in \mathcal B'$ has a single, well-defined $\bigl(A_m(b), \kappa_{\pole(m)}(b)\bigr)$,
so as $(a,\kappa)$ ranges over all valid pairs, each $b$ contributes on the right-hand side
exactly once, at its own pair, and then once more summed over $k$; by $(\mathrm{Disagg})$
that inner sum over $k$ equals $\mu_{b,t}$. Hence the right-hand side is exactly
$\sum_{b \in \mathcal B'} \mu_{b,t}$, and the identity above reads $1 = \sum_b \mu_{b,t}$.
Had $(\mathrm{Cons})$ been imposed only for $a=2$, the ``not maintained'' part of $q$ would
be left unconstrained by any column, and nothing would force this sum to $1$.
\end{remark}

\begin{remark}
Dualising $(\mathrm{Cons})$ gives, for a fixed stage $t$, the dual vector
$\lambda_t^{m,a,k,\kappa} := \mathrm{dual}\bigl((\mathrm{Cons})_{m,a,k,\kappa,t}\bigr)$
(\texttt{consistency\_duals}), used by the pricing problem of Section~\ref{sec:pricing}.
The dual of $(\mathrm{Disagg})$ is never used: that constraint is created afresh with $b$
itself, so it carries no information available to future pricing iterations.
\end{remark}

\section{Pricing MILP of the Full Fluid LP}
\label{sec:pricing}

Coded in \texttt{fluid\_lp\_cg.jl}, function \texttt{solve\_pricing}. Solved once per
stage $t \in [H]$ at every column-generation iteration, given the dual vector
$\lambda_t = (\lambda_t^{m,a,k,\kappa})$ of Section~\ref{sec:master}. This is problem
(5)/(13) of the paper: find the column of smallest reduced cost
$\breve c(b,t) = f(b) + \sum_{m} \lambda_t^{m,A_m(b),k,\kappa_{\pole(m)}(b)}$, where $k$ is
now also chosen by the pricing problem (see Remark~\ref{rem:pricing-k}).

\subsection*{Variables}
\begin{align*}
  z_m &\in \{0,1\} &&\forall m \in [M]
    &&\text{($m$ is maintained this period)}\\
  s_{m,\tau} &\in \{0,1\} &&\forall m \in [M],\ \tau \in [L_m]
    &&\text{(maintenance of $m$ starts on day $\tau$; not created if $L_m<1$, then $z_m{=}0$)}\\
  \any_{p,\tau} &\in [0,1] &&\forall p \in [K],\ \tau \in [T]
    &&\text{(relaxed: some component of $p$ under maintenance on day $\tau$)}\\
  \free_{p,\tau} &\in \{0,1\} &&\forall p \in [K],\ \tau \in [T]
    &&\text{(day $\tau$ of pole $p$ declared penalty-free)}\\
  w_{p,\tau} &\in [0,1] &&\forall p \in [K],\ \tau \in [T]
    &&\text{(relaxed linearisation of $\any_{p,\tau} \wedge \free_{p,\tau}$)}\\
  \ksel_{p,(k,\kappa)} &\in \{0,1\} &&\forall p \in [K],\ 0 \le \kappa \le k \le K_0(p)
    &&\text{(pole $p$'s quota is $k$ and $\kappa$ free days are used)}\\
  y_{m,k,\kappa} &\in [0,1] &&\forall m \in [M],\ 0 \le \kappa \le k \le K_0(\pole(m))
    &&\text{(linearisation of $z_m \wedge \ksel_{\pole(m),(k,\kappa)}$).}
\end{align*}
\begin{remark}[Pricing must also choose $k$]
\label{rem:pricing-k}
Since the master now disaggregates $\mu_{b,t}$ by the pole's remaining quota $k$
(Section~\ref{sec:master}), pricing must decide, for every pole, not only how many free
days $\kappa$ to use but also \emph{at which quota level $k \ge \kappa$} the new mass is
attached: this is why the selection variable $\ksel_{p,\kappa}$ of a first version is
enlarged from $\kappa$ alone to valid pairs $(k,\kappa)$. The McCormick linearisation
$y_{m,k,\kappa} = z_m \wedge \ksel_{\pole(m),(k,\kappa)}$ has the exact same structure as
before, just re-indexed.

\textcolor{red}{The $k$ found by the pricing MILP is only a \emph{witness}: it certifies
that, for pole $p$, there exists at least one quota level $k \ge \kappa$ at which schedule
$b$ has negative reduced cost, which is enough to justify adding $b$ to $\mathcal B'$. Once
$b$ is accepted, \texttt{add\_column!} does \emph{not} restrict itself to that particular
$k$: it creates $\mu^{(p)}_{k,b,t}$ for \emph{every} valid $k \in \{\kappa_p(b),\dots,K_0(p)\}$
(Section~\ref{sec:master}). It is then the master LP itself that decides, across all of
these, how much mass to actually place at each $k$.}
\end{remark}
\begin{remark}
$\any$ and $w$ are declared \emph{continuous} in $[0,1]$ (not binary) in the code: since
the objective below penalises $\any_{p,\tau}$ and rewards $w_{p,\tau}$, and both are
sandwiched between binary quantities by the constraints below, they take integral values
at every optimal solution --- relaxing them only speeds up the solver.
\end{remark}

\begin{remark}[Only two of the three usual AND-linearisation constraints]
\label{rem:w-linearisation}
$w_{p,\tau}$ linearises $\any_{p,\tau} \wedge \free_{p,\tau}$, and the textbook
linearisation of an AND of two binaries needs three constraints:
$w \le \any$, $w \le \free$, and $w \ge \any + \free - 1$ (to stop $w$ from being picked
\emph{smaller} than the true AND). Below, only the first two are imposed. This is safe
here because $w_{p,\tau}$ enters the objective solely through the term
$-\mathrm{pen}_p(\tau)\, w_{p,\tau}$, with $\mathrm{pen}_p(\tau) \ge 0$, in a
\emph{minimisation}: increasing $w_{p,\tau}$, up to its upper bound
$\min(\any_{p,\tau}, \free_{p,\tau})$, only ever decreases the objective and is never
penalised elsewhere, so the solver always pushes $w_{p,\tau}$ up to exactly
$\min(\any_{p,\tau}, \free_{p,\tau})$ at any optimum --- precisely the AND value once
$\any_{p,\tau}, \free_{p,\tau} \in \{0,1\}$. The third constraint would therefore never
change the optimal solution; it is dropped as redundant, matching the pricing MILP of
\texttt{column\_generation/column\_generation.jl} that this code was adapted from. (The
Compact Fluid LP of Section~\ref{sec:compact} keeps all three inequalities on its own
$w_{p,\tau,t}$, by the same reasoning applying just as well there; they are kept there
only for extra caution, not because they are needed.)
\end{remark}
Define the auxiliary expression (not a variable; \texttt{ong} in the code)
\[
  \ong_{m,\tau} \;:=\; \sum_{s\,=\,\max(1,\tau-d_m+1)}^{\min(\tau,L_m)} s_{m,s}
  \qquad(\text{$1$ iff $m$ is under maintenance on day }\tau).
\]

\subsection*{Constraints}
\begin{align}
  \sum_{\tau=1}^{L_m} s_{m,\tau} &= z_m
    &&\forall m \in [M] \tag{10}\\
  \sum_{m=1}^{M} \rho_m\, \ong_{m,\tau} &\le R
    &&\forall \tau \in [T] \tag{11}\\
  \any_{p,\tau} &\ge \ong_{m,\tau}
    &&\forall p \in [K],\ m \in \mathrm{comps\_of\_pole}(p),\ \tau \in [T] \\
  w_{p,\tau} &\le \free_{p,\tau}, \qquad w_{p,\tau} \le \any_{p,\tau}
    &&\forall p \in [K],\ \tau \in [T] \\
  \sum_{0 \le \kappa \le k \le K_0(p)} \ksel_{p,(k,\kappa)} &= 1
    &&\forall p \in [K] \tag{12a}\\
  \sum_{\tau=1}^{T} \free_{p,\tau} &= \sum_{0 \le \kappa \le k \le K_0(p)} \kappa\, \ksel_{p,(k,\kappa)}
    &&\forall p \in [K] \tag{12b}\\
  y_{m,k,\kappa} &\le z_m, \qquad y_{m,k,\kappa} \le \ksel_{\pole(m),(k,\kappa)}, \qquad
    y_{m,k,\kappa} \ge z_m + \ksel_{\pole(m),(k,\kappa)} - 1
    &&\forall m \in [M],\ \kappa \le k \le K_0(\pole(m))
\end{align}

\subsection*{Objective (reduced cost of Eq.~(13))}
\[
  \min\;\; \underbrace{\sum_{p=1}^K \sum_{\tau=1}^T \mathrm{pen}_p(\tau)\,
      \bigl(\any_{p,\tau} - w_{p,\tau}\bigr)}_{=\ f(b)}
  \;+\;
  \sum_{m=1}^{M}\ \sum_{0 \le \kappa \le k \le K_0(\pole(m))}
    \Bigl[\, \lambda_t^{m,1,k,\kappa}\, \ksel_{\pole(m),(k,\kappa)}
    \;+\; \bigl(\lambda_t^{m,2,k,\kappa} - \lambda_t^{m,1,k,\kappa}\bigr)\, y_{m,k,\kappa} \,\Bigr].
\]
\begin{remark}
The bracketed term equals $\sum_m \lambda_t^{m,A_m(b),k^\star,\kappa^\star}$: for a
component with $z_m=0$ (not maintained) $y_{m,k,\kappa}=0$ for every $(k,\kappa)$, so its
contribution reduces to $\lambda_t^{m,1,k^\star,\kappa^\star}$ where $(k^\star,\kappa^\star)$
is the pair selected by $\ksel_{\pole(m),\cdot}$; for a component with $z_m=1$, the
McCormick constraints force $y_{m,k,\kappa}=\ksel_{\pole(m),(k,\kappa)}$, so the
contribution reduces to $\lambda_t^{m,2,k^\star,\kappa^\star}$. Since $\ksel_{p,\cdot}$ is
shared by every component of pole $p$, all of them agree on the same $(k^\star,\kappa^\star)$
--- consistently with the master's variable $\mu^{(p)}_{k^\star,b,t}$ being common to the
whole pole.
\end{remark}

\subsection*{Extraction of the column}
$\mathrm{start}_m = \tau$ if $z_m=1$ and $s_{m,\tau}=1$, else $0$; $\mathrm{free}_p =
\{\tau : \free_{p,\tau}=1\}$. The column is added to $\mathcal B'$ (Section
\ref{sec:master}) whenever its reduced cost is below a tolerance $-\mathrm{tol}$.

\section{Light Fluid LP: quota as a subsystem}
\label{sec:light}

Coded in \texttt{fluid\_cg\_light.jl}, functions \texttt{build\_master\_light} /
\texttt{add\_column\_light!} (master LP) and \texttt{solve\_pricing\_light} (pricing MILP).
This is an alternative to the Master LP of Section~\ref{sec:master}: instead of giving every
component $m$ of a pole its own copy of the pole's quota state, so that $q$ carries
$(x,k,\kappa)$ jointly (Section~\ref{sec:master}'s local state $(x,k)$ and local action
$(a,\kappa)$), we introduce one extra \emph{quota subsystem} per pole $p$, whose local state
is $k \in \{0,\dots,K_0(p)\}$ and local action is $\kappa \in \{0,\dots,k\}$, and let every
real component $m$ carry only its own $(x,a)$. This drops the exact correlation between a
component's state and its pole's quota level that Section~\ref{sec:master}'s design tracks,
giving a (weakly) looser bound $Z^{\mathrm{light}} \le Z$, but removes the $(k,\kappa)$
blow-up of every component's index and the disaggregated $\mu^{(p)}_{k,b,t}$ variables
entirely: consistency ties each real component straight to $\mu_{b,t}$ on its own mode, and
each quota subsystem straight to $\mu_{b,t}$ on its own $\kappa$, with no per-$k$
bookkeeping needed on either side.

\subsection*{Variables}
\begin{align*}
  q_{m,x,a,t} &\ge 0
  &&\forall m \in [M],\ x \in [n_m],\ a \in \{1,2\},\ t \in [H], \\
  q^{\kappa}_{p,k,\kappa,t} &\ge 0
  &&\forall p \in [K],\ 0 \le \kappa \le k \le K_0(p),\ t \in [H], \\
  \mu_{b,t} &\ge 0
  &&\forall b \in \mathcal B' \subseteq \mathcal B,\ t \in [H],
\end{align*}
where $q^{\kappa}_{p,k,\kappa,t}$ is the probability that pole $p$'s quota subsystem is in
state $k$ and takes action $\kappa$ at stage $t$ (\texttt{qk} in the code); $q_{m,x,a,t}$ no
longer carries $(k,\kappa)$ at all, unlike Section~\ref{sec:master}.

\subsection*{Objective}
\[
  \min \; \sum_{m,x,a,t} c_m(x,a)\, q_{m,x,a,t}
  \;+\; \sum_{b \in \mathcal B',\,t} f(b)\, \mu_{b,t},
\]
with $c_m$ and $f$ exactly as in Section~\ref{sec:master} ($q^{\kappa}$ carries no cost).

\subsection*{Constraints}

\textbf{Initial condition} ($x_0, k_0$ the current observed state):
\[
  \sum_{a=1}^{2} q_{m,x,a,1} \;=\; \mathbb 1\{x = x_{0,m}\}
  \qquad \forall m \in [M],\ x \in [n_m],
\]
\[
  \sum_{\kappa=0}^{k} q^{\kappa}_{p,k,\kappa,1} \;=\; \mathbb 1\{k = k_{0}(p)\}
  \qquad \forall p \in [K],\ k \in \{0,\dots,K_0(p)\}.
\]

\textbf{Flow of a real component} ($t \ge 2$; ordinary single-component MDP flow, with no
quota entanglement at all):
\[
  \sum_{a=1}^{2} q_{m,x,a,t}
  \;=\;
  \sum_{x' \in [n_m]}\ \sum_{a=1}^{2} p^a_m(x \mid x')\; q_{m,x',a,t-1}
  \qquad \forall m \in [M],\ x \in [n_m]. \tag{Flow-$q$}
\]

\textbf{Flow of a quota subsystem, no reset between $t{-}1$ and $t$} ($t \ge 2$):
\[
  \sum_{\kappa=0}^{k} q^{\kappa}_{p,k,\kappa,t}
  \;=\;
  \sum_{\kappa'=0}^{K_0(p)-k} q^{\kappa}_{p,\,k+\kappa',\,\kappa',\,t-1}
  \qquad \forall p \in [K],\ k \in \{0,\dots,K_0(p)\}. \tag{Flow-$q^\kappa$}
\]

\textbf{Flow of a quota subsystem, reset between $t{-}1$ and $t$} ($t \ge 2$):
\[
  \sum_{\kappa=0}^{k} q^{\kappa}_{p,k,\kappa,t} \;=\; 0 \quad \text{if } k \ne K_0(p),
  \qquad
  \sum_{\kappa=0}^{K_0(p)} q^{\kappa}_{p,K_0(p),\kappa,t}
  \;=\;
  \sum_{k'=0}^{K_0(p)}\ \sum_{\kappa'=0}^{k'} q^{\kappa}_{p,k',\kappa',t-1}
  \quad \text{if } k = K_0(p). \tag{Flow-$q^\kappa$-reset}
\]

\textbf{Consistency, real components:}
\[
  \sum_{x \in [n_m]} q_{m,x,a,t}
  \;-\;
  \sum_{\substack{b \in \mathcal B':\\ A_m(b)=a}} \mu_{b,t}
  \;=\; 0
  \qquad \forall m \in [M],\ a \in \{1,2\},\ t \in [H]. \tag{Cons-$q$}
\]

\textbf{Consistency, quota subsystems:}
\[
  \sum_{k=\kappa}^{K_0(p)} q^{\kappa}_{p,k,\kappa,t}
  \;-\;
  \sum_{\substack{b \in \mathcal B':\\ \kappa_p(b)=\kappa}} \mu_{b,t}
  \;=\; 0
  \qquad \forall p \in [K],\ \kappa \in \{0,\dots,K_0(p)\},\ t \in [H]. \tag{Cons-$q^\kappa$}
\]

Every new column $b$ is added (\texttt{add\_column\_light!}) by, for every $t$: creating
$\mu_{b,t}$ (objective coefficient $f(b)$) and entering it with coefficient $-1$ in
$(\mathrm{Cons}\text{-}q)_{m,A_m(b),t}$ for every $m \in [M]$, and in
$(\mathrm{Cons}\text{-}q^\kappa)_{p,\kappa_p(b),t}$ for every $p \in [K]$ -- no
disaggregation step is needed, unlike Section~\ref{sec:master}.

\textbf{Non-negativity:} $q, q^\kappa, \mu \ge 0$.

\subsection*{Pricing MILP}

Coded in \texttt{fluid\_cg\_light.jl}, function \texttt{solve\_pricing\_light}. Dualising
$(\mathrm{Cons}\text{-}q)$ and $(\mathrm{Cons}\text{-}q^\kappa)$ gives, at a fixed stage $t$,
duals $\lambda_t^{m,a} := \mathrm{dual}((\mathrm{Cons}\text{-}q)_{m,a,t})$ and
$\lambda^{\kappa}_{t,p,\kappa} := \mathrm{dual}((\mathrm{Cons}\text{-}q^\kappa)_{p,\kappa,t})$
(\texttt{consistency\_duals\_light}). The pricing problem is
\[
  \min_{b \in \mathcal B} \;\; f(b) \;+\; \sum_{m=1}^M \lambda_t^{m,A_m(b)}
  \;+\; \sum_{p=1}^K \lambda^{\kappa}_{t,p,\kappa_p(b)}.
\]
Since neither dual depends jointly on $(a,\kappa)$ any more, the reduced cost is already
linear in the schedule variables separately: it uses the same variables
$z_m, s_{m,\tau}, \any_{p,\tau}, \free_{p,\tau}, w_{p,\tau}$ as Section~\ref{sec:pricing}, and
a selector $\kappa\mathrm{sel}_{p,\kappa} \in \{0,1\}$, $\kappa \in \{0,\dots,K_0(p)\}$
(\emph{without} $k$, unlike $\ksel_{p,(k,\kappa)}$ of Section~\ref{sec:pricing}), subject to
constraints (10), (11) and
\[
  \sum_{\kappa=0}^{K_0(p)} \kappa\mathrm{sel}_{p,\kappa} = 1,
  \qquad
  \sum_{\tau=1}^{T} \free_{p,\tau} = \sum_{\kappa=0}^{K_0(p)} \kappa\, \kappa\mathrm{sel}_{p,\kappa}
  \qquad \forall p \in [K],
\]
with objective
\[
  \min\;\; \underbrace{\sum_{p=1}^K \sum_{\tau=1}^T \mathrm{pen}_p(\tau)\,
      \bigl(\any_{p,\tau} - w_{p,\tau}\bigr)}_{=\ f(b)}
  \;+\; \sum_{m=1}^M \Bigl[\lambda_t^{m,1} + \bigl(\lambda_t^{m,2}-\lambda_t^{m,1}\bigr) z_m\Bigr]
  \;+\; \sum_{p=1}^K \sum_{\kappa=0}^{K_0(p)} \lambda^{\kappa}_{t,p,\kappa}\, \kappa\mathrm{sel}_{p,\kappa}.
\]
\begin{remark}
No McCormick linearisation is needed: the component term is already linear in $z_m$ (it
equals $\lambda_t^{m,1}$ at $z_m=0$ and $\lambda_t^{m,2}$ at $z_m=1$), and the quota term is
already linear in $\kappa\mathrm{sel}_{p,\cdot}$, since $\lambda^\kappa_{t,p,\kappa}$ does not
depend on $k$. This is the direct counterpart of Remark~\ref{rem:pricing-k}: there is no $k$
left for the pricing problem to choose a witness for, since $q^\kappa$ is not disaggregated
by $b$ in the first place.
\end{remark}

\section{Order-2 higher-order formulation (Section 4.2)}
\label{sec:order2}

Coded in \texttt{fluid\_cg\_order2.jl}, function \texttt{build\_master\_order2}: adds
pairwise pseudomarginals on top of the master LP of Section~\ref{sec:master} (built
unchanged by \texttt{build\_master}), without touching the pricing MILP of
Section~\ref{sec:pricing} or the columns $\mathcal B'$ at all.

\subsection*{Why pairwise pseudomarginals help here}
The failure cost $\delta_m(x)$ of Section~\ref{sec:master} is computed component by
component, assuming every \emph{other} component of the pole is at full capacity
(\texttt{failure\_cost} in \texttt{instances.jl}). This is exactly the source of
looseness of Section 4.1: the true capacity of a pole is the \emph{minimum} of its
components' individual capacities (\texttt{pole\_penalty}), not additively separable, so
summing $\delta_m$ over the components of a pole under-estimates the cost of two of them
being degraded \emph{simultaneously}. Cross-pole pairs never interact this way: the
penalty decomposes additively across poles (\texttt{instance\_penalty}), so for
$m_1, m_2$ of \emph{different} poles, summing $\delta_{m_1}(x_1) + \delta_{m_2}(x_2)$ is
already \emph{exact} --- see the remark below for how this is handled.

\subsection*{The pairwise correction $\Delta$}
For two components $m_1 \ne m_2$ of the \emph{same} pole $p$ (with
$c_i := \mathrm{comp\_of}(m_i)$, $P_{\mathrm{evac}}$ and $P_{\mathrm{daily}}$ as in
\texttt{instances.jl}), define the joint failure cost of the pair when neither is
maintained, propagating their (independent, given no maintenance) daily degradation
forward from $(x_1,x_2)$:
\[
  \delta^{\mathrm{joint}}_{m_1,m_2}(x_1,x_2) \;:=\; \sum_{\tau=1}^{T}
  \mathrm{pen}_p(\tau) \left(1 - \frac{1}{100} \sum_{s_1,s_2}
  v_1^{(\tau)}(s_1)\, v_2^{(\tau)}(s_2)\, \min\bigl(P_{\mathrm{evac}}(c_1,s_1),\,
  P_{\mathrm{evac}}(c_2,s_2)\bigr) \right),
\]
where $v_i^{(1)} := \mathbb 1_{x_i}$ and $v_i^{(\tau+1)} := P_{\mathrm{daily}}(c_i)\,
v_i^{(\tau)}$ (the marginal degradation of component $m_i$ evolved day by day, weather-free).
The pairwise correction is
\[
  \Delta_{m_1,m_2}(x_1,x_2) \;:=\; \delta^{\mathrm{joint}}_{m_1,m_2}(x_1,x_2)
  - \delta_{m_1}(x_1) - \delta_{m_2}(x_2) \;\le\; 0 \qquad (m_1, m_2 \text{ of the same pole}),
\]
(never positive: capacity under the true minimum-of-two model is always at least as good
as under the additive model, so the joint cost is never larger than the sum of the two
separate costs). For $m_1, m_2$ of \emph{different} poles, $\delta^{\mathrm{joint}}$ has no
physical meaning (there is no shared capacity to take a minimum of), and instead
$\Delta_{m_1,m_2} \equiv 0$ is set directly, by the exact-decomposition argument above.

\begin{remark}[Which pairs are tracked]
By default, \texttt{compute\_Δ\_pairs} computes $\Delta_{m_1,m_2}$ and tracks \emph{every}
pair of distinct components, same-pole \emph{and} cross-pole alike --- all
$\binom{M}{2}$ of them (e.g.\ $\binom{7}{2}=21$ for one pole, $\binom{14}{2}=91$ for two).
Cross-pole pairs are legitimate to request and are never rejected; since
$\Delta_{m_1,m_2} \equiv 0$ for them, \texttt{build\_master\_order2} recognises this and
skips building their (expensive) variables and constraints entirely, so tracking them
costs nothing beyond computing the (trivially zero) correction once. Two restriction
mechanisms exist but are \emph{off} unless explicitly requested, in the same spirit as
lift-and-project (Section~\ref{sec:lift-project}, ``adding only a few well chosen cuts''):
if \texttt{max\_pair\_states} is given, pairs with $n_{m_1} n_{m_2} >
\texttt{max\_pair\_states}$ (the most expensive to track) are dropped first; and if
\texttt{max\_pairs} is given, only that many of the remaining pairs, ranked by largest
impact $\max_{x_1,x_2} |\Delta_{m_1,m_2}(x_1,x_2)|$, are kept (a pair whose
$\Delta$ correction is negligible everywhere is not worth its cost). Since every
cross-pole pair has impact exactly $0$, this ranking naturally sorts them last, so
restricting to a small \texttt{max\_pairs} automatically prioritises same-pole pairs
first, without needing to treat the two kinds of pairs separately. Both keywords are
exposed by \texttt{make\_policy\_fluid\_cg\_order2} for instances where tracking every
pair is not tractable; a user-supplied explicit list \texttt{pairs} bypasses both filters
entirely and may freely mix same-pole and cross-pole pairs.
\end{remark}

\subsection*{Variables}
For every pair $(m_1,m_2)$ with $\Delta_{m_1,m_2} \not\equiv 0$ (i.e.\ every \emph{same-pole}
pair that was kept; cross-pole pairs are skipped entirely, see the remark above), writing
$p$ for their common pole:
\[
  p^{\mathrm{pair}}_{m_1,m_2,x_1,x_2,a_1,a_2,k,\kappa,t} \ge 0
  \qquad \forall x_1 \in [n_{m_1}],\ x_2 \in [n_{m_2}],\ a_1,a_2 \in \{1,2\},\
  0 \le \kappa \le k \le K_0(p),\ t \in [H]
\]
(both local modes $a_1, a_2$, not only ``neither maintained''; see
Remark~\ref{rem:order2-both-modes} below for why both are needed).

\subsection*{Constraints}
\textbf{Initial condition} (product of the two deterministic initial states, mirroring
(Init)):
\[
  \sum_{a_1,a_2,\kappa} p^{\mathrm{pair}}_{m_1,m_2,x_1,x_2,a_1,a_2,k_0(p),\kappa,1}
  \;=\; \mathbb 1\{x_1=x_{0,m_1}\} \cdot \mathbb 1\{x_2=x_{0,m_2}\}
  \qquad \forall x_1, x_2.
\]

\textbf{Flow} (mirroring (Flow)/(Flow-reset), with the \emph{product} transition kernel
$p_{m_1}^{a_1}(x_1\mid x_1') \, p_{m_2}^{a_2}(x_2\mid x_2')$, since the two components
evolve independently given their own local actions --- Definition~1 of the paper): for
every $t \ge 2$, $x_1, x_2$, $k_n \in \{0,\dots,K_0(p)\}$, writing
$\sigma_n := \sum_{a_1,a_2,\kappa \le k_n} p^{\mathrm{pair}}_{m_1,m_2,x_1,x_2,a_1,a_2,k_n,\kappa,t}$,
\[
  \sigma_n = 0 \ \text{ if reset and } k_n \ne K_0(p); \qquad \text{otherwise}
\]
\[
  \sigma_n \;=\; \sum_{x_1',x_2'}\ \sum_{a_1,a_2}\ p_{m_1}^{a_1}(x_1\mid x_1')\,
  p_{m_2}^{a_2}(x_2\mid x_2') \!\!
  \sum_{\substack{(k,\kappa) \to k_n \\ \text{(same predecessors as (Flow))}}} \!\!
  p^{\mathrm{pair}}_{m_1,m_2,x_1',x_2',a_1,a_2,k,\kappa,t-1}.
\]

\textbf{Marginal consistency with $q$, at every stage} (not only $t=1$):
\[
  \sum_{x_2,a_2} p^{\mathrm{pair}}_{m_1,m_2,x_1,x_2,a_1,a_2,k,\kappa,t} = q_{m_1,x_1,a_1,k,\kappa,t}
  \quad \forall x_1, a_1, \qquad
  \sum_{x_1,a_1} p^{\mathrm{pair}}_{m_1,m_2,x_1,x_2,a_1,a_2,k,\kappa,t} = q_{m_2,x_2,a_2,k,\kappa,t}
  \quad \forall x_2, a_2.
\]
This is the general definition of a \emph{consistent} order-$2$ pseudomarginal (Section
4.2: ``marginalising a higher-order pseudomarginal down to a lower order recovers it''),
specialised from constraint (9c) of Proposition~9's LP (Appendix A of the paper) to a pair
of \emph{distinct} local actions $a_1,a_2$ instead of one action shared by both.

\subsection*{Objective}
Only the ``neither maintained'' slice $a_1=a_2=1$ carries a $\Delta$ term (it is the only
mode combination for which $\delta^{\mathrm{joint}}_{m_1,m_2}$ was computed):
\[
  \sum_{(m_1,m_2)} \ \sum_{k,\kappa,t} \ \sum_{x_1,x_2}
  \Delta_{m_1,m_2}(x_1,x_2)\; p^{\mathrm{pair}}_{m_1,m_2,x_1,x_2,1,1,k,\kappa,t}
\]
added to the objective of Section~\ref{sec:master}.

\begin{remark}[Why $p^{\mathrm{pair}}$ needs both modes, and both Flow and marginal consistency]
\label{rem:order2-both-modes}
A first version of this master defined $p^{\mathrm{pair}}$ only for $a_1=a_2=1$, tied to
$q$ solely by the two inequalities $\sum_{x_2} p^{\mathrm{pair}} \le q_{m_1,x_1,1,k,\kappa,t}$
and symmetrically for $m_2$, with no Flow constraint at all. A second version added Flow
(the product-kernel recursion above) but dropped the marginal-consistency constraint,
reasoning that it followed from Flow by induction on $t$. \emph{That argument is wrong}:
$q$'s own (Flow) constraint (Section~\ref{sec:master}) only pins the sum over \emph{both}
mode \emph{and} $\kappa$ at every stage, never an individual mode --- so nothing forces
$p^{\mathrm{pair}}$'s mode split to track $q$'s actual (column-determined) mode split
beyond $t=1$, and even at $t=1$ only the sum over $(a_1,a_2,\kappa)$ is pinned, not the
individual split. Without an explicit tie, the solver is free to route as much mass as it
likes into the $a_1=a_2=1$ slice (rewarded by $\Delta \le 0$) independently of what $q$ is
actually doing, which is unsound: in testing it let the order-2 bound fall \emph{below}
the order-1 one, the reverse of a valid tightening. Both Flow (which pins the
\emph{correlation} between $x_1$ and $x_2$ via genuine dynamics, avoiding an
adversarially-chosen joint shape) and marginal consistency at every $t$ (which pins the
mode split to $q$'s real one) are needed together, mirroring how Proposition~9's own LP
(Appendix A) uses both (9c) (consistency) and (9e) (a flow-like recursion) for its
pairwise variable.
\end{remark}

\begin{remark}
$p^{\mathrm{pair}}$ constrains $q$ only through equalities that do not touch $\mu$: no
column of $\mathcal B'$ is priced against it, so the pricing MILP of
Section~\ref{sec:pricing} and the duals $\lambda_t^{m,a,k,\kappa}$ of
Section~\ref{sec:master} are used completely unchanged --- column generation for this
order-2 master proceeds exactly as in Algorithm~1, just against a larger master LP.
\end{remark}

\section{Compact Fluid LP}
\label{sec:compact}

Coded in \texttt{compact\_lp.jl}, function \texttt{build\_compact\_LP}. This is Eq.~(6) of
the paper (Section 3.3), specialised as in Appendix C.4: the integrality of the joint
action polytope is relaxed and one continuous copy of it is embedded \emph{per period},
which removes the need for column generation at the price of a looser bound. It reuses
exactly the $q_{m,x,a,k,\kappa,t}$ variables and the (Init), (Flow), (Flow-reset)
constraints of Section~\ref{sec:master} --- \emph{unchanged} --- and replaces
$(\mathrm{Cons})$/$(\mathrm{Disagg})$ by a relaxed, per-period copy of the pricing schedule
of Section~\ref{sec:pricing}.

\begin{remark}[No $(k,b)$ strengthening here]
Since the Compact Fluid LP has no discrete action $b$ at all (it is replaced by a
continuous relaxation of the schedule polytope, one copy per period), the strengthening of
Remark~\ref{rem:strengthen} does not apply: there is no $\mu^{(p)}_{k,b,t}$ to share
between components. We therefore keep, unchanged, the original \emph{intra-pole quota
copy} constraint --- every component of a pole shares the same $(k,\kappa)$ marginal as a
reference component $m_1(p) := \mathrm{comps\_of\_pole}(p)[1]$:
\[
  \sum_{x \in [n_m]}\ \sum_{a=1}^{2} q_{m,x,a,k,\kappa,t}
  \;=\;
  \sum_{x \in [n_{m_1(p)}]}\ \sum_{a=1}^{2} q_{m_1(p),x,a,k,\kappa,t}
  \tag{Quota-copy}
\]
for all $p \in [K]$, $m \in \mathrm{comps\_of\_pole}(p) \setminus \{m_1(p)\}$,
$0 \le \kappa \le k \le K_0(p)$, $t \in [H]$.
\end{remark}

\begin{remark}[Warm start across states: one template per $(t_0^{\mathrm{cyc}}, H)$, not one
model per state]
\label{rem:compact-warmstart}
Unlike the Full Fluid LP, whose master is warm-started across states through its reused
\emph{columns} (Section~\ref{sec:master}), the Compact Fluid LP has no such combinatorial
object to carry over: its structure (every variable, every constraint, every coefficient)
depends only on $(t_0^{\mathrm{cyc}}, H)$, never on the specific state $(x_0, k_0)$ being
solved for -- $x_0$ and $k_0$ only ever enter through (a) the right-hand side of the
initial-condition constraint (Init) and (b) which $q_{m,x,a,k,\kappa,1}$ are fixed to $0$ at
$t=1$ (only the $k_0(p)$ slice of each pole is free). \texttt{build\_compact\_template}
(\texttt{compact\_lp.jl}) therefore builds the model once per $(t_0^{\mathrm{cyc}}, H)$ --
in practice a handful of distinct pairs, since $H$ only shrinks near the very end of the
$n_{\mathrm{years}}$-year horizon and $t_0^{\mathrm{cyc}} \in [6]$ -- with every $q_{\cdot,1}$
fixed to $0$ and every (Init) right-hand side at $0$. \texttt{configure\_compact\_LP!} then
points that SAME model (mutated in place, never rebuilt) at each new $(x_0,k_0)$ sharing that
$(t_0^{\mathrm{cyc}}, H)$: it updates the $M \times \max_m n_m$ (Init) right-hand sides via
\texttt{set\_normalized\_rhs}, and re-fixes/unfixes only the (at most) two $k$-slices that
change per pole (the previous $k_0(p)$ back to $0$, the new one to free) rather than
touching every $k$. Because only bounds and right-hand sides change -- no variable or
constraint is ever added or removed -- Gurobi's simplex (\texttt{Method=1}, dual simplex, in
place of the barrier method without crossover used previously, which has no comparable warm
start) resumes from the previous solve's basis instead of solving from scratch.
\texttt{make\_policy\_compact\_lp} keeps one template per $(t_0^{\mathrm{cyc}}, H)$ for the
whole simulation and reconfigures it for every new state, instead of calling
\texttt{build\_compact\_LP} (still available, e.g.\ for a one-off solve, but now just
\texttt{build\_compact\_template} immediately followed by a single
\texttt{configure\_compact\_LP!}) for every one. Measured on \texttt{one\_pole\_simple}
(MTBFs and technician shares of Section~\ref{sec:simplified}), this reduced the mean time
per solve from $0.21$s to $0.02$s (a $10\times$ speed-up), making \texttt{compact} faster
per solve than \texttt{fluid\_cg} ($0.07$s) on the same instance -- consistent with
\texttt{compact} needing no pricing MILP at all, unlike \texttt{fluid\_cg}.
\end{remark}

\subsection*{Additional variables (one full copy of the schedule per stage $t \in [H]$)}
\begin{align*}
  z_{m,t} &\in [0,1], \quad
  s_{m,\tau,t} \in [0,1]\ (\tau \in [L_m]), \quad
  \any_{p,\tau,t} \in [0,1], \\
  \free_{p,\tau,t} &\in [0,1], \quad
  w_{p,\tau,t} \in [0,1], \quad
  \ksel_{p,\kappa,t} \in [0,1]\ (\kappa \le \bar K,\ \text{fixed }0\text{ if } \kappa>K_0(p)).
\end{align*}
with $\ong_{m,\tau,t} := \sum_{s=\max(1,\tau-d_m+1)}^{\min(\tau,L_m)} s_{m,s,t}$.

\subsection*{Objective}
\[
  \min\;\; \sum_{m,x,a,k,\kappa,t} c_m(x,a)\, q_{m,x,a,k,\kappa,t}
  \;+\; \sum_{t=1}^{H}\sum_{p=1}^{K}\sum_{\tau=1}^{T} \mathrm{pen}_p(\tau)\,
        \bigl(\any_{p,\tau,t} - w_{p,\tau,t}\bigr).
\]

\subsection*{Constraints}
(Init), (Flow), (Flow-reset) as in Section~\ref{sec:master}, (Quota-copy) as restated
above, plus, for every $t \in [H]$:
\begin{align}
  \sum_{\tau=1}^{L_m} s_{m,\tau,t} &= z_{m,t}
    &&\forall m \in [M] \tag{10$'$}\\
  \sum_{m=1}^{M} \rho_m\, \ong_{m,\tau,t} &\le R
    &&\forall \tau \in [T] \tag{11$'$}\\
  \any_{p,\tau,t} &\ge \ong_{m,\tau,t}
    &&\forall p,\ m \in \mathrm{comps\_of\_pole}(p),\ \tau \\
  w_{p,\tau,t} &\le \free_{p,\tau,t}, \quad
  w_{p,\tau,t} \le \any_{p,\tau,t}, \quad
  w_{p,\tau,t} \ge \free_{p,\tau,t} + \any_{p,\tau,t} - 1
    &&\forall p,\ \tau \\
  \sum_{\kappa=0}^{\bar K} \ksel_{p,\kappa,t} &= 1, \qquad
  \sum_{\tau=1}^{T} \free_{p,\tau,t} = \sum_{\kappa=0}^{\bar K} \kappa\, \ksel_{p,\kappa,t}
    &&\forall p \tag{12$'$}
\end{align}
and the two coupling constraints tying the relaxed schedule to the state-action variables
$q$ (Appendix C.4):
\[
  \sum_{x \in [n_m]}\ \sum_{k=0}^{K_0(\pole(m))}\ \sum_{\kappa=0}^{k} q_{m,x,2,k,\kappa,t}
  \;\le\; z_{m,t}
  \qquad \forall m \in [M],\ t \in [H], \tag{Coupling-$z$}
\]
\[
  \sum_{x \in [n_{m_1(p)}]}\ \sum_{a=1}^{2}\ \sum_{k=0}^{K_0(p)}\ \sum_{\kappa=0}^{k}
    \kappa\; q_{m_1(p),x,a,k,\kappa,t}
  \;=\; \sum_{\tau=1}^{T} \free_{p,\tau,t}
  \qquad \forall p \in [K],\ t \in [H]. \tag{Coupling-$\kappa$}
\]

\begin{remark}[A slack exploited by the relaxation]
\label{rem:slack}
$(\mathrm{Coupling}\text{-}z)$ is an \emph{inequality}, exactly as written in Appendix C.4
of the paper, so nothing directly forces $z_{m,t}$ down to the actual replacement mass
$\sum_{x,k,\kappa} q_{m,x,2,k,\kappa,t}$. This alone would not be a problem if inflating
$z_{m,t}$ always had a cost elsewhere (a legitimate optimum would then keep it tight): a
larger $z_{m,t}$ forces more $m_{m,\tau,t}'$ mass (Eq.~$(10')$), hence more
$\ong_{m,\tau,t}$ and $\any_{p,\tau,t}$, which \emph{does} cost $\mathrm{pen}_p(\tau)$
per day \emph{unless} that day is also declared free ($w_{p,\tau,t}=\any_{p,\tau,t}$).
The reason this cost can always be dodged is a second, compounding slack:
$(\mathrm{Coupling}\text{-}\kappa)$ sums $\kappa \, q_{m_1(p),x,a,k,\kappa,t}$ over
\emph{both} modes $a$, not only $a=2$. Since the mode-$1$ (``not maintained'') mass
completely dominates $q$ in practice (components rarely fail), and nothing else in the
LP ties a component's $\kappa$-attribution while unmaintained to anything physical, the
solver is free to book a sizeable share of that mode-$1$ mass at $\kappa > 0$, manufacturing
$\free_{p,\tau,t}$ supply through $(\mathrm{Coupling}\text{-}\kappa)$ out of thin air. That
manufactured free-day supply is then exactly what pays for the inflated, purely fictitious
schedule of $z_{m,t}$, closing the loop at zero net objective cost. This was confirmed
numerically: for a component whose true replacement mass is $O(10^{-11})$ (so $a=2$ is
essentially unused) and $z_{m,1} \approx 0.21$, the mode-$1$ mass at $t=1$ splits as
$\approx 0.61$ at $\kappa=0$ and $\approx 0.39$ at $\kappa=1$ --- the model is ``spending''
a free day on a component it has no intention of maintaining, purely to subsidise
$z_{m,t}$ elsewhere. Consequently, the policy-extraction step
(\texttt{extract\_policy\_compact}) reads the maintenance decision from the replacement
mass $\sum_{x,k,\kappa} q_{m,x,2,k,\kappa,1}$ rather than from $z_{m,1}$.
\end{remark}

\subsection*{Policy extraction (rounding and repair heuristic)}
The compact LP is solved once per re-optimisation (no pricing loop). Its solution is
generally fractional and, per Section 3.5 of the paper, is turned into an admissible
action of stage $t=1$ by a repair heuristic (not itself part of the LP, and not to be
confused with the LP variables it reads from): components with replacement mass
$\sum_{x,k,\kappa} q_{m,x,2,k,\kappa,1}$ above $\tfrac12$ are the candidates, examined in
decreasing order of $\rho_m$ (their own technician requirement of constraint $(11)$ ---
placing the components hardest to fit, resource-wise, first). For each candidate $m$, the
relaxed start-day weights $s_{m,\tau,1}$, $\tau \in \{1,\dots,L_m\}$, of Eq.~$(10')$ are
used to rank candidate days \emph{twice over}: first among the days already opened by an
earlier candidate of the same pole (highest $s_{m,\tau,1}$ first), to concentrate
maintenance on as few distinct days as possible --- essential for the free-day quota to
cover them all, see the discussion after Eq.~$(12')$ --- and only if none of those fits
$(11')$, among every remaining day in $\{1,\dots,L_m\}$ (again by decreasing
$s_{m,\tau,1}$); a component that fits nowhere is left unmaintained this period. Free days
are then chosen, per pole, among the days that end up under maintenance, capped at the
quota $k_p$ actually remaining and prioritising the highest-penalty days.

\section{Lift-and-project cuts (Section 4.3)}
\label{sec:lift-project}

Coded in \texttt{lift\_project.jl}. Proposition 10 of the paper bounds
$q^{m_1}_{x_1,a,t} + q^{m_2}_{x_2,a,t}$ by $\mu_{a,t}$ plus a correction term
$\tau^+_{m_1,m_2,t}(x_1,x_2)$ obtained by \emph{propagating}, through one step of the
transition kernel, a bound on the order-1 marginals of stage $t-1$; this is well suited to
a D-MDP where every subsystem shares the exact same global action $a$ (Definition 1), so
that $q^{m_1}_{x_1,a,t}$ and $q^{m_2}_{x_2,a,t}$ can both be written against that same $a$.

\subsection*{Why this needs to be adapted}
In our setting there is no such shared global action: a joint action (column) $b$ assigns
independent local modes $A_{m_1}(b) \ne A_{m_2}(b)$ in general. Propagating a
transition-based bound as in Proposition 10 is not the right tool here, because our
columns already give us something the D-MDP toy example does not have: for every
generated column $b \in \mathcal B'$, we know \emph{exactly} which local mode it assigns
to every component, so the joint probability that ``$m_1$ is maintained this way \emph{and}
$m_2$ is maintained that way'' is directly and exactly computable from the existing
$\mu_{b,t}$ variables, with no need to go through order-1 information from a previous
stage at all. The cut below uses this instead.

\subsection*{A first attempt, and why it is very loose on this problem}
\label{sec:lift-project-first-attempt}
A first version of this note bounded $\mathbb P[E \cup F]$ by
$\nu_{p,(k,\kappa),t} := \sum_{x,a} q_{m,x,a,k,\kappa,t}$ (any component $m$ of pole $p$,
Remark~\ref{rem:sum-to-one}) and $\mathbb P[E \cap F]$ by the \emph{exact} probability
$\rho_{m_1,m_2,(a_1,a_2,\kappa),t} := \sum_{b \in \mathcal B': A_{m_1}(b)=a_1,\,
A_{m_2}(b)=a_2,\, \kappa_p(b)=\kappa} \mu_{b,t}$ that a column maintaining $m_1$ \emph{and}
$m_2$ in exactly the prescribed way was chosen. Both quantities are exact, but this
combination is a poor fit for the offshore substation instance:
$\nu_{p,(k,\kappa),t}$ ignores $a_1, a_2$ entirely (it aggregates over \emph{every} mode of
\emph{every} component of the pole), so it is far looser than necessary, while
$\rho_{m_1,m_2,(a_1,a_2,\kappa),t}$ requires a \emph{single} generated column to maintain
$m_1$ and $m_2$ \emph{simultaneously} in the exact prescribed way, which technician
capacity (11) makes rare: with $R$ technicians shared by up to $M=7$ components,
few (if any) generated columns maintain two specific components on the very same day in
the very same period, so $\rho$ is $0$, or close to it, for most $(a_1,a_2)$, making the
cut needlessly tight on the wrong term while staying needlessly loose on the other.

\subsection*{The adapted cut}
We instead bound each side of the inclusion--exclusion identity with the tool best suited
to it: the \emph{union} $E \cup F$, which only asks that \emph{at least one} of the two
prescribed local actions was used, is bounded exactly from the already-tracked
disaggregated variables $\mu^{(p)}_{k,b,t}$ of Section~\ref{sec:master}; the
\emph{intersection} $E \cap F$, which asks about the joint \emph{state} of $m_1$ and
$m_2$ regardless of their actions, is bounded as in Proposition 10, by propagating a
$\min$ of order-1 marginals of stage $t-1$ through the transition kernel.

Fix a pole $p$, two distinct components $m_1 \ne m_2 \in \mathrm{comps\_of\_pole}(p)$,
modes $a_1, a_2 \in \{1,2\}$, a valid quota pair $0 \le \kappa \le k \le K_0(p)$, and a
stage $t \ge 2$ (a transition is propagated, so, as in Proposition 10, $t=1$ is excluded --
harmlessly, since the state is deterministic at $t=1$ anyway). Define
\[
  \sigma_{m_1,m_2,(a_1,a_2,k,\kappa),t} \;:=\;
  \sum_{\substack{b \in \mathcal B':\\ A_{m_1}(b)=a_1 \ \text{or}\ A_{m_2}(b)=a_2,\ \ \kappa_p(b)=\kappa}}
  \mu^{(p)}_{k,b,t},
\]
\[
  \tau^+_{m_1,m_2,t}(x_1,x_2;k) \;:=\;
  \sum_{(k',\kappa') \to k} \ \sum_{a_1',a_2' \in \{1,2\}} \ \sum_{x_1',x_2'}
  p_{m_1}^{a_1'}(x_1 \mid x_1')\, p_{m_2}^{a_2'}(x_2 \mid x_2')\,
  \min\bigl(q_{m_1,x_1',a_1',k',\kappa',t-1},\, q_{m_2,x_2',a_2',k',\kappa',t-1}\bigr),
\]
where ``$(k',\kappa') \to k$'' ranges over the same predecessor quota states summed by
(Flow)/(Flow-reset) for target $k$ at stage $t$ (Section~\ref{sec:master}), and $a_1', a_2'$
are dummy (summed-out) mode labels for stage $t-1$, unrelated to the fixed $a_1, a_2$ of
the cut itself. Let
$E := \{X_{m_1,t}=x_1,\ A_{m_1}(B_t)=a_1,\ K_{p,t}=k,\ \mathrm K_{p,t}=\kappa\}$ and
$F := \{X_{m_2,t}=x_2,\ A_{m_2}(B_t)=a_2,\ K_{p,t}=k,\ \mathrm K_{p,t}=\kappa\}$, so
$\mathbb P[E] = q_{m_1,x_1,a_1,k,\kappa,t}$ and $\mathbb P[F] = q_{m_2,x_2,a_2,k,\kappa,t}$
exactly, as before.
\begin{itemize}
\item $E \cup F$ implies $\{(A_{m_1}(B_t)=a_1 \text{ or } A_{m_2}(B_t)=a_2),\ K_{p,t}=k,\
  \mathrm K_{p,t}=\kappa\}$ (dropping the state conditions $X_{m_1,t}=x_1$, $X_{m_2,t}=x_2$
  only enlarges the event), whose probability is \emph{exactly}
  $\sigma_{m_1,m_2,(a_1,a_2,k,\kappa),t}$ (a disjoint union over the qualifying columns),
  so $\mathbb P[E \cup F] \le \sigma_{m_1,m_2,(a_1,a_2,k,\kappa),t}$.
\item $E \cap F$ implies $\{X_{m_1,t}=x_1,\ X_{m_2,t}=x_2,\ K_{p,t}=k\}$ (dropping the
  action conditions $A_{m_1}(B_t)=a_1$, $A_{m_2}(B_t)=a_2$, $\mathrm K_{p,t}=\kappa$ this
  time), whose probability is bounded, exactly as $\tau_{m_1,m_2,t}(x_1,x_2)$ in the proof
  of Proposition 10, by propagating the transition kernel one step from stage $t-1$:
  $\mathbb P[E \cap F] \le \tau^+_{m_1,m_2,t}(x_1,x_2;k)$.
\end{itemize}
By inclusion--exclusion,
\[
  q_{m_1,x_1,a_1,k,\kappa,t} + q_{m_2,x_2,a_2,k,\kappa,t}
  \;=\; \mathbb P[E] + \mathbb P[F]
  \;=\; \mathbb P[E \cup F] + \mathbb P[E \cap F]
  \;\le\; \sigma_{m_1,m_2,(a_1,a_2,k,\kappa),t} + \tau^+_{m_1,m_2,t}(x_1,x_2;k),
  \tag{C$'$}
\]
which holds for every occupation measure realisable by an admissible policy, for all
$p \in [K]$, $m_1 \ne m_2 \in \mathrm{comps\_of\_pole}(p)$, $x_1 \in [n_{m_1}]$,
$x_2 \in [n_{m_2}]$, $a_1, a_2 \in \{1,2\}$, $0 \le \kappa \le k \le K_0(p)$, $t \ge 2$.

\begin{remark}
$\sigma$ is tighter than the first attempt's $\nu$ because it only counts columns
consistent with \emph{at least one} of the two prescribed actions, using the
already-tracked $\mu^{(p)}_{k,b,t}$ rather than the coarser, action-blind $q$-aggregate.
$\tau^+$ replaces the first attempt's $\rho$ because it asks a strictly easier question of
the generated columns: whether $m_1$ and $m_2$ can \emph{jointly reach} states
$(x_1,x_2)$ from \emph{any} pair of order-1 histories at $t-1$ (via the same
$\min$-of-marginals argument as Proposition 10), rather than whether a single column was
ever generated that maintains both of them in the exact prescribed way at $t$ -- the
latter being needlessly hard to satisfy under a shared technician-capacity constraint,
which is exactly the source of the "very suboptimal" cut of the first attempt.
\end{remark}

\section{Policy extraction: a method-agnostic heuristic}
\label{sec:heuristic-extraction}

Coded in \texttt{heuristic\_extraction.jl} (function \texttt{extract\_policy\_heuristic},
\texttt{schedule\_true\_optimal}). Every \texttt{make\_policy\_*} function of
Sections~\ref{sec:master}--\ref{sec:lift-project} takes a keyword argument
\texttt{heuristic\_extraction}, \texttt{true} by default: instead of letting each method's
own machinery also decide the \emph{schedule} (which day each maintained component starts,
which specific days are declared free) -- the column with the largest weight
$\mu_{b,1}$ for fluid\_cg/fluid\_cg\_light/fluid\_cg\_order2/lift\_project, or the
LP-weight-based greedy day picker of \texttt{extract\_policy\_compact} for the Compact Fluid
LP -- every method now goes through the \emph{same} two-step rule.

\begin{remark}[Why the schedule needs a shared rule]
The additive per-component cost used by every formulation above,
$c_m(x,a) = \delta_m(x)$ if $a=1$ and $\mathrm{cc}_m$ if $a=2$ (Section~\ref{sec:notation}),
never depends on \emph{which day} of the period a maintained component ($a=2$) actually
starts: replacing a component is charged $\mathrm{cc}_m$ flat, as if the repair happened
instantly, with no term charging the days spent still-degraded while waiting for the
repair to start. The true, non-additive cost (\texttt{pole\_penalty}, \texttt{instances.jl})
has no such blind spot: a component left in a bad state for longer keeps contributing to
the pole's capacity loss every day until it is actually repaired. Since the additive
proxy is indifferent between an early and a late start of the very same maintained
component, which day gets picked is entirely a tie-break of the underlying solver (column
pool composition and weights, or LP degeneracy) -- and nothing stops it from landing on
the worst one. This was observed directly in \texttt{results/report\_one\_pole.txt}:
scenario 100, period 119, \texttt{fluid\_cg\_order2} started the repair of a
badly-degraded component on day 41 instead of day 1 (both give the exact same
$\mathtt{lp\_cost} = 0.5025$), incurring $40$ extra units of true penalty from $40$ days of
avoidable pre-repair exposure that \texttt{compact} (which happened to tie-break toward day
1) did not pay.
\end{remark}

\paragraph{Step 1 -- read the aggregate decision off $q$.} For every component $m$, at the
known current state $x_0(m)$, read
\[
  \mathrm{mode}_m \;:=\; \argmax_{a \in \{1,2\}}\ q_{m,x_0(m),a,K_0(\pole(m)),\kappa,1}
  \ \text{summed over}\ \kappa \in \mathrm{qrange}(K_0(\pole(m))).
\]
For every pole $p$, at the known current quota $K_0(p)$, read
\[
  \kappa_p \;:=\; \argmax_{\kappa \in \mathrm{qrange}(K_0(p))}\
  \sum_{a=1}^{2} q_{m_1(p),x_0(m_1(p)),a,K_0(p),\kappa,1},
\]
using the pole's reference component $m_1(p)$ (every component of a pole shares the same
$(k,\kappa)$ distribution, Section~\ref{sec:notation}, so any component would give the same
answer). For \texttt{fluid\_cg\_light} (Section~\ref{sec:light}), whose $q$ has no $(k,\kappa)$
index, the same two reads use $q_{m,x_0(m),a,1}$ directly for $\mathrm{mode}_m$ and the
pole's own quota-subsystem variable $q^\kappa_{p,K_0(p),\kappa,1}$ for $\kappa_p$.

\paragraph{Step 2 -- re-derive the schedule for the true cost.} Given the maintained set
$S = \{m : \mathrm{mode}_m = 2\}$ and $(\kappa_p)_{p\in[K]}$, function
\texttt{schedule\_true\_optimal} builds the actual $(\mathrm{start}_m)_{m\in S}$ and free
days by:
\begin{itemize}
  \item scheduling every $m \in S$ on the \emph{earliest} feasible day of the start-day grid
  ($1, 1+\mathrm{QS}, 1+2\,\mathrm{QS}, \dots$, Section~\ref{sec:pricing}) compatible with
  the shared technician capacity $R$, most urgent component first (components are sorted
  by decreasing $\delta_m(x_0(m))$, ties broken by decreasing $\rho_m$, and packed by a
  simple earliest-fit greedy against a single shared demand profile, since $R$ is shared
  across every pole). Since $\delta_m(x) \ge 0$ and repairing a component only ever resets
  it to the as-new state, starting earlier is always weakly better for the true cost --
  never worse -- so this is not a further approximation on top of the $S$ chosen at Step 1,
  it removes exactly the blind spot described above. A component the greedy cannot fit this
  period (a rare artefact of Step 1 deciding $\mathrm{mode}_m$ independently for every
  component) is simply left unmaintained rather than left infeasible;
  \item placing the $\kappa_p$ free days of pole $p$ on the first $\kappa_p$ days (in day
  order) among the days some component of $p$ is actually under maintenance -- the only
  days \texttt{pole\_penalty} ever waives on a declared free day (a free day with no
  maintenance ongoing changes nothing: \texttt{instances.jl}, lines 109--117).
\end{itemize}

Passing \texttt{heuristic\_extraction = false} restores each method's original, bespoke
extraction (column argmax alone, or \texttt{extract\_policy\_compact}), kept for comparison.
The command-line driver \texttt{run\_experiment.jl} exposes this as \texttt{--no-heuristic}.

\section{Baseline 1: two-step optimisation (simplified problem)}
\label{sec:baseline1}

Coded in \texttt{baseline1.jl} (\texttt{optimize\_thresholds}, \texttt{schedule\_min\_cost},
\texttt{make\_policy\_baseline1}), run with \texttt{run\_experiment.jl baseline1
one\_pole\_simple ...}. It is only defined for the simplified instances of
Section~\ref{sec:simplified} (it raises an error otherwise) and uses no LP relaxation.

\paragraph{Step 1: one threshold policy per component (offline, once per instance).}
Each component $m$ is optimised alone, among the policies ``maintain iff $x \ge \theta_m$'',
$\theta_m \in \{1,\dots,n_m+1\}$ ($\theta_m = n_m+1$ means never maintain). Component $m$ alone
is a Markov chain over the periods: in state $x$,
\[
  \text{if } x < \theta_m:\ \text{cost } \delta_m(x),\ x \to x' \sim P^{\mathrm{period}}_m(\cdot \mid x);
  \qquad
  \text{if } x \ge \theta_m:\ \text{cost } \delta_m(x) + \mathrm{cc}_m + \gamma\, d_m\,\bar{\mathrm{pen}}_p,\ x \to 1,
\]
where $\bar{\mathrm{pen}}_p = \frac1T\sum_\tau \mathrm{pen}_p(\tau)$ ($=1$ here), so the
maintenance cost is $\gamma$ times the duration $d_m$ of the maintenance, in addition to the
replacement cost $\mathrm{cc}_m$. The failure cost $\delta_m(x)$ of the period is paid in both
cases (Section~\ref{sec:simplified}): maintaining does not avoid the failure cost of the
period in which it happens, it only resets the state for the next periods, so a preventive
maintenance can pay off.
$\gamma \in [0,1]$ is a parameter of the method (\texttt{gamma}, \texttt{--gamma=<value>} in
\texttt{run\_experiment.jl}, default $\gamma = 1$; a value outside $[0,1]$ is rejected). It
represents the fact that several components are often maintained at the same time and share
the same maintenance days, so the whole cost of these days cannot be charged to a single
component: $\gamma < 1$ charges only a fraction of them to $m$, and $\gamma = 1$ charges all
of them. $\gamma$ only enters this step: Step~2 below always uses the true costs.
\textcolor{red}{In theory $\gamma$ could depend on the component (a component that is
usually maintained together with others would get a smaller $\gamma_m$ than one that is
usually maintained alone). For now, the same $\gamma$ is used for every component: the
policies are computed with a single $\gamma$, independent of the component.}
$\theta_m$ is the
value minimising the expected total cost of this chain over the horizon of the simulation
($n_{\mathrm{years}} \times 6$ periods, starting in state $1$), computed exactly by propagating
the state distribution (ties go to the smallest $\theta_m$). This step ignores the
technician-capacity constraint and any coupling between components or poles.
$\gamma$ only acts through the thresholds, which take a finite number of values as $\gamma$
varies: the larger $\gamma$, the later a component is maintained. \texttt{tune\_gamma.jl}
scans a fine grid of $\gamma \in [0,1]$, keeps the distinct threshold vectors and simulates
each of them on the same $1000$ scenarios. On \texttt{one\_pole\_simple} (MTBFs of the
original instances) the best vector is $\theta = (11,2,4,3,3,3,3)$, obtained for
$\gamma \in [0.944, 1]$ (mean cost $56.6$, standard error $1.5$); the next ones are $57.7$
for $\gamma \in [0.568, 0.922]$ and $60.0$ for $\gamma \le 0.555$. (For reference, a scan up
to $\gamma = 100$, before restricting $\gamma$ to $[0,1]$, found nothing better: larger
$\gamma$ only give later thresholds and higher costs, from $63.1$ up to $644.9$.) The
default $\gamma = 1$ lies in the best range, so it is kept.

\paragraph{Step 2: schedule for the joint state (online).} In the simulation, at the joint
state $x$, the set of components to maintain is $S = \{m : x_m \ge \theta_m\}$. The schedule
is then the one that maintains exactly the components of $S$, with the smallest period cost.
\begin{itemize}
  \item If a schedule maintaining all of $S$ exists (start days on the grid, lead margin,
  shared technician capacity (11)), the schedule is the one with the fewest maintenance days
  $f(b)$ among those: since the set of maintained components is fixed, this is the minimiser
  of the period cost $\sum_m \delta_m(x_m) + \sum_{m \in \mathrm{maintained}(b)} \mathrm{cc}_m + f(b)$
  (Section~\ref{sec:simplified}).
  \item Otherwise, the components of $S$ are considered by decreasing priority
  $x_m - \theta_m$ (how far the state is above the threshold; ties: decreasing $\delta_m(x_m)$).
  Each one is kept if a feasible schedule still exists for the set kept so far plus this
  component, and is otherwise left unmaintained for this period (it will be considered again
  at the next period, as long as $x_m \ge \theta_m$). The schedule is then the one with the
  fewest maintenance days for the final set.
\end{itemize}
Feasibility and the schedule are computed with the pricing MILP of
Section~\ref{sec:pricing}, in which the maintenance of the wanted set is imposed by a large
negative cost on ``maintained'' (and a huge cost for the components outside the set); the
remaining objective is then $f(b)$ plus the replacement costs, constant for a given set.
The schedule is applied for the period, and the procedure is repeated at the next period.
Schedules are cached by joint state.
\end{document}
